\section{Magnetic trial actions with unequal walls}
\label{sec:07_magnetic_analysis-trials}

The next estimates retain a positive part of the real energy while gaining
$2L_X$ from the common interval and controlling the activities on the two
shoulders. We first define the trial action, then separate its smoothing,
common-interval, and shoulder estimates. The resulting bounds will be
implemented in the actual particle integral in the next section.

\subsection{Explicit particles and the trial action}

Retain the magnetic contours, their elementary multiplicities \(w_I\),
and the real-energy form \(B\) of the preceding section. A double label
always represents the sum of its two elementary constituents when a kernel
is computed. Write $s_X=s$, $s_Y=b$, $\bar\nu_X=\nu_X$ and
$\bar\nu_Y=\rho$ for their sizes and positive unit-mass mean densities.
The two main labels are \(X,Y\); the subscript \(n\)
means the pair of main labels. Throughout this section,
\[
 \beta=\frac\pi4,\qquad \gamma=\frac\pi{3\beta}=\frac43,
 \qquad \kappa=(3\beta)^2,\qquad
 p=(3-\theta_\infty)\beta,\qquad L=L_Y.
\]
The parameter \(p\) in the kernel formulas is distinct from a marked port.
We assume \(c_1L\le L_X\le L_Y=L\), with \(c_1>0\) fixed.
Constants can depend on \(c_1\). An error \(o(L)\), or \(o(1)B(x)\),
is uniform over the stated frozen configurations after all smoothing,
localization, and projection parameters have been fixed. When counts are
restricted to a fixed power of \(L\), uniformity includes the positions
and counts in that restriction. Errors denoted \(o_\varepsilon(1)L\)
have a coefficient tending to zero as \(\varepsilon\downarrow0\), with
\(h=\varepsilon^{K_2}\) as specified below; they are different from the
large-\(L\) errors.

Split main particles at their respective walls \([-L_g,L_g]\).
Put the particles outside their own walls, together with every far-label
particle, into the explicit configuration \(\zeta\). Its total number is
\(n_e\), with elementary multiplicities included whenever a group total
is tested. The interior density \(\alpha_g\) has mass \(m'_g\), obtained
by subtracting these explicit multiplicities from \(s-M,b+M\).
Here \(M\) remains an integer. We use the following explicit regularization. Fix a small $h>0$.
Let $\psi=\psi^{(h)}$ be the even positive convolution probability
whose Fourier transform averages
$\cosh((\pi-hv)k)/\cosh(\pi k)$ over $0<v<1$, with probability
density proportional to $e^{-1/\sqrt v}$. Before averaging its
real density is
\[
 \frac{\cosh(u/2)\sin(hv/2)}
 {2\pi(\sinh^2(u/2)+\sin^2(hv/2))}.
\]
Hence $\psi\le C/h$, its tails are at most $Ch e^{-|u|/2}$ for
$|u|\ge1$, and $\int u^2\psi(u)\,du\le Ch$. Its Fourier
transform has upper and lower stretched-exponential bounds of the
form $C e^{-c(h|k|)^{1/3}}$ and $c e^{-C(h|k|)^{1/3}}$.
Put $S_0(u)=-2\log|2\sinh(u/2)|$. Then
$S_0\ge S_0*\psi\ge S_0*\psi*\psi$ by the preceding Fourier
formula with the real tails fixed. Choose an even Gevrey cutoff
$\chi$, equal to one on $[-1/2,1/2]$, zero outside $[-1,1]$, and
between zero and one. On like-main interior/explicit entries set
\[
 R_h(u)=\chi(u)(S_0-S_0*\psi)(u)\ge0,\qquad
 \widetilde H_{ne}=H_{ne}-R_h,
\]
with $R_h=0$ on other entries. We also write $R$ for this kernel
when no operator complement is involved. This cutoff uses the
horizontal difference, including on the curved contour. Define
\[
 \mathcal E(\alpha,\zeta)
 =\tfrac12(\alpha-\mu,H_{nn}(\alpha-\mu))
   +(\alpha-\mu,\widetilde H_{ne}\zeta)
   +\tfrac12(\zeta,H\zeta)_{\rm off}.
\]
The last expression omits each individual explicit self interaction;
near empirical diagonals will instead use smooth reference compensation. Put
\[
 x=\operatorname{Re}\alpha-\mu+\psi*\zeta,\qquad
 z=\operatorname{Im}\alpha,\qquad
 D'=\sum_Iw_Ix_I,\qquad D_z'=\sum_Iw_Iz_I.
\]
The primitives vanish at infinity, and each main component of \(z\) has
integral zero. To distinguish a contour from a density, write
\(\mathfrak z_I(u)\) for its complex parametrization; thus
\(\mathfrak z_X(u)=u+i\beta\theta(u)\) on main \(X\).
Let \(\mathfrak t_I\) be the sum of the particle logarithms from the two
marked anomalies. The trial action is
\[
 \mathcal D(\alpha,\zeta)
 =\mathcal E(\alpha,\zeta)
   -(\mathfrak t_n,\alpha-\mu)-(\mathfrak t,\zeta).
\]
A purely imaginary constant depending only on the species counts never
affects a bound on its real part.

\subsection{Reference measure and sector normalization}\label{sec:reference-law}

The comparison law has two independent labels, denoted $X,Y$. For an
integer $n\ge0$, its unnormalized measure on $[0,\pi]^n$ is
\[
 \frac1{n!}\prod_{i=1}^n\frac{2\sin t_i\,dt_i}{2\pi}
       \prod_{i<j}|2(\cos t_i-\cos t_j)|^2.
\]
Under $x=\cos t$ this is the Legendre orthogonal-polynomial ensemble.
Andr\'eief's determinant integration identity gives its mass
\cite[equation~(1.3)]{TracyWidom1998}
\begin{equation}\label{eq:reference-normalization}
 J_n=\prod_{i=0}^{n-1}
 \frac{2\cdot16^i(i!)^4}{\pi(2i+1)((2i)!)^2}
 =\exp\{-\tfrac14\log(n+1)+O(1)\},\qquad J_0=1.
\end{equation}
We write $\mathbb E_0$ only for the probability expectation, divided
by $J_{N_X}J_{N_Y}$. Neither of these factors is part of the
observable.

On each label, let $r\in[-A_g,A_g]$ be a real reference parameter
and let $t_g(r)$ be an increasing probability coordinate taking values
in $[0,\pi]$, with density $\lambda_g(r)=t'_g(r)/\pi$.
The corresponding point on the main contour is denoted by $Z_g(r)$.
We will specify this parametrization, including its endpoint caps, in
Section~\ref{sec:07_magnetic_analysis-fluctuations}. With $X_{N_g}$
the empirical counting measure in the reference parameter, put
\[
 q_g=X_{N_g}-N_g\lambda_g(r)\,dr,\qquad Q_g'=q_g,
 \qquad Q_g(-A_g)=Q_g(A_g)=0.
\]
The reference kernel is block diagonal, with each diagonal block
\[
 H_{\rm ref}(t,t')=-2\log|2(\cos t-\cos t')|.
\]
Its mean against $dt/\pi$ is zero. The Hilbert space $\mathcal X$
is the completion of smooth Dirichlet primitives for the form
\[
 h_0(Q,Q)=(Q',H_{\rm ref}Q')
    =\pi^2\sum_{g,n\ge1}n|b_{g,n}|^2,
 \qquad Q_g(t)=\sum_{n\ge1}b_{g,n}\sin(nt).
\]
Real symmetric forms extend bilinearly to complex arguments; norms
and adjoints always use the Hermitian convention. An empirical
primitive need not have finite $h_0$ energy. We evaluate its quadratic
cost only after subtracting the reference singularity, so that the
remaining kernel is smooth, including its diagonal. No subtraction of
two infinite self energies is involved.

Let $p_g$ be the pushforward of $\lambda_g(r)\,dr$ to the horizontal
coordinate of the $g$ wall, and write $p_L=(p_X,p_Y)$. For this
sector take $N_g=m'_g$, restricted to nonnegative integer counts.
Inserting $m'p_L+q$ into the trial action always means pushing $q$
to the same horizontal coordinates. The two interior empirical
measures are denoted by $X_{\rm int}=(X_{N_X},X_{N_Y})$; their
pushforward is likewise understood in pairings with the trial kernels.

We now compute the factor produced by changing to this reference law.
Put $a_*=2\cos(\beta/2)=d/c$. On a main label the local particle
factor and the reference pair factor have the respective expansions
\[
\begin{aligned}
 E_s(Z_g(r)-Z_g(r'))
   &=a_*^2Z'_g(r)^2(r-r')^2(1+o(1)),\\
 |2(\cos t_g(r)-\cos t_g(r'))|^2
   &=\bigl(2\sin t_g(r)t'_g(r)\bigr)^2(r-r')^2(1+o(1)).
\end{aligned}
\]
Consequently the smooth-diagonal compensation contributes
\[
 \exp\!\left(\frac{(H-H_{\rm ref})_{gg}(r,r)}2\right)
 =\frac{2\sin t_g(r)t'_g(r)}{a_*Z'_g(r)}.
\]
Here $H$ is composed with the contour parametrization, and the branch
is continued from the straight contour; endpoint values are taken by
continuity. Multiplication by $dZ_g(r)/(2\pi)$ cancels $Z'_g(r)$
and leaves $a_*^{-1}$ times the reference one-particle measure
$2\sin t_g\,dt_g/(2\pi)$.

The finite identity and ordinary mean extraction contribute
$((2d)C_*)^n$, where $n=N-2$ is the total elementary particle degree.
Since
\[
 \frac{(2d)C_*}{a_*}=\frac\gamma2,
 \qquad
 N_X+N_Y=n-\sum_{v\in\zeta}w_{I(v)},
\]
extracting $(\gamma/2)^n$ leaves $a_*^{w_I}$ for each explicit
particle of species $I$. Thus a singleton retains activity $d/c$,
and a double retains $4c^2(d/c)^2$. The magnetic mean identities
supply all ordinary and marked site-pair factors. Restoring the
missing marked--marked factor and replacing $n$ by $N$ in the
extracted power changes the result only by a uniformly bounded
factor, because the allowed marked log differences are $0,\pm i\beta$
and $H_*$ is nonzero there.

After extracting $(\gamma/2)^N\prod_{i<j}H_*(a_i-a_j)^2$, the
remaining sector is therefore an unnormalized reference integral
with integrand
\[
 A\exp\!\left(-\mathcal D(m'p_L+q,\zeta)
                       +\tfrac12H_{\rm ref}[q,q]\right),
 \qquad A=\exp(-(X_{\rm int},R\zeta)).
\]
The difference in the exponent uses smooth-diagonal compensation on
both labels, as above. Including the explicit integration measures,
the remaining prefactor is bounded in modulus by
\[
 C\prod_{v\in\zeta}
 \left(\frac{|d\mathfrak z_{I(v)}(v)|}{2\pi}
 (d/c)^{w_{I(v)}} a_{I(v)}\,e^{(\mu,R\delta_v)}\right),
\]
with a factorial divisor for each explicit species. Here $a_I=1$ for
a main or far singleton and $a_I=4c^2$ for a double; on straight
contours $d\mathfrak z=dv$. Count-dependent signs have modulus one.
Finally $(\mu,R\delta_v)=O_h(1)$, and is $o(1)$ past a macroscopic
buffer outside the corresponding wall.


\subsection{Smoothing and count control}

Smoothing lets the positive real energy control the explicit particles.
The first estimate counts them locally; the second controls their total
number, group imbalance, and distance beyond the longer wall.

\begin{lemma}[Curved smoothing and count control]\label{lem:07_magnetic_analysis-smoothing}
Let \(n_j\) be the local explicit counts on the smoothing grids, with a separate grid beginning at each wall. Uniformly for sufficiently small \(\varepsilon>0\),
\begin{equation}
 \sum n_j^2\le C(1+|\log h|)(B(x)+1),\qquad
 \ell_e+(M^2+n_e^2)/L\le C_h (B(x)+1),
 \quad\ell_e=\sum_{v\in\zeta}(|v|-L_Y)_+ .
 \label{eq:M-20}
\end{equation}


Put
\[
 \Delta_H=\tfrac12\{(\zeta,\Re H\zeta)_{\rm off}
                      -(\psi*\zeta,\Re H\,\psi*\zeta)\}.
\]
Then
\begin{equation}
\begin{aligned}
 \Re\mathcal E={}&B(x)-B(z)-(x,\Im H z)+\Delta_H
     \\ &+O\!\left(\epsilon_{\rm sm}(h,\varepsilon,L)(B(x)+B(z)+1)\right),\\
 \epsilon_{\rm sm}(h,\varepsilon,L)\le{}& C\varepsilon^{-C}h^{c_2}+o_L(1).
\end{aligned}
 \label{eq:M-21}
\end{equation}
Here \(C,c_2>0\); \(o_L(1)\to0\) at fixed \(h,\varepsilon\). Decreasing \(c_2\) absorbs logarithmic losses in \(h\). Also, entrywise off the omitted diagonals,
\begin{equation}
 \Re(H-\psi H\psi)_{IJ}(u,v)
 \ \ge\ -C\varepsilon^{-C}h^{c_2}e^{-c_2|u-v|}
            -o_L(1)e^{-c_2|u-v|}\quad\text{at explicit positions}.
 \label{eq:M-22}
\end{equation}
Here \(\psi H\psi\) denotes separate convolution. Its own-variable diagonal costs at most \(C(1+|\log h|)\) (upper bound on real part, uniform in \(\varepsilon\)).


\end{lemma}
\begin{proof}
\par\smallskip\noindent\textbf{Positive count tests.}
We give the count tests explicitly. For each far label use translates
of a nonnegative bump covering a unit cell to test the positive field
$x_I=\psi*\zeta_I$. For a main label use a grid beginning at its own
wall and bumps supported outside that wall. The first bump rises from
zero to one across distance $h$, remains one through distance two,
and then cuts off. Every particle in that first cell contributes a
fixed positive amount after smoothing. The squared $H^{1/2}$ norm
of a bump is $O(1+|\log h|)$, as follows from the difference-quotient
seminorm, and bounded overlaps give the same bound for linear
combinations with the sum of squared coefficients. The mean cell
masses are summably bounded because $s_g\bar\nu_g(u)=O(1)$ at
its own wall. Duality and magnetic coercivity therefore give the
first count estimate.

Beyond the larger wall, a count to the right of $v$ is bounded by
twice its positive smeared cumulative there. Its square integrates
to $O(B(x)+1)$ using the total primitive and the exponentially small
mean tail; the left end is identical. Counts are integers, so this
also bounds their distance sum $\ell_e$. The remaining count mass
lies in $O(L)$ cells and is bounded by Cauchy--Schwarz. Finally the
$X$-group total of $x$ is $-M$. Test it on a smooth plateau covering
the larger wall interval, whose $H^{1/2}$ norm is $O(\sqrt L)$;
its complement is bounded by the positive smeared tail masses.
This proves the remaining count estimates. The covariance kernel
of $(1+k^2)^{-1/2}$ is bounded by
$C(1+|\log|u||)e^{-c|u|}$; convolving twice with $\psi$ also gives
$\|\psi*\zeta\|_{-1/2}^2\le C(1+|\log h|)\sum_j n_j^2$,
which we use in the smoothing estimate below.

\par\smallskip\noindent\textbf{Horizontal singularity and curved residuals.}
We give the additional kernel estimates needed for curvature. Subtract \(S_0\delta_{IJ}\) from \(H_{IJ}\). This accounts for the full horizontal angle-zero terms, using double fusion; the residual curved self-log ratio on main has bounded size locally by divided differences and the graph estimates. Except for this ratio, all residual log singularities have horizontal offsets in imaginary height bounded away from zero by \(c_3\varepsilon\), after double fusion, for like near-diagonal real positions. This follows from \(1/2\le\theta\le1\), the singletons \(25/8,-7/2\), doubles \(3+\varepsilon,4+\varepsilon;-3,-4\), and the shifts in \(E_s,E_c\); for two curved main particles with different real positions use uniform Lipschitzness near their diagonal. All angle-nonzero sums thus have fixed-order local derivative bounds polynomial in \(1/\varepsilon\) (up to one derivative in each argument suffices, or fixed higher order where the curve has regular derivatives). In the far tails of the difference coordinate use the convergent log series. More precisely, the remaining leading terms after subtracting \(S_0\delta_{IJ}\) are
\[
 -(w_Iw_J-\delta_{IJ})\log(2\cosh(u-v))
 +i(c^0_{IJ}-\upsilon_I(u) w_J+w_I \upsilon_J(v))\tanh(u-v)
\]
with exponentially localized residuals. Here \(c^0_{IJ}\) are straight-contour sign coefficients, \(\upsilon_X(u)=\beta(\theta(u)-\theta_\infty)\) on main \(X\), and \(\upsilon_I=0\) otherwise. Indeed vary just the imaginary coordinate there: each reduced-angle log contributes its height translation to the linear tail, including the continued squared diagonal log via the same log ratio, and there is no zero crossed off a self diagonal. Derivatives up to one in each variable of the exponentially localized bounds also hold away from the curved same-main diagonal neighborhood. In that neighborhood the extra log ratio is smooth with uniform fixed-derivative bounds outside a fixed neighborhood of zero; near zero it is bounded with first derivative in \(u\) at most \(C/(|u|+d_0)\) when the graph approaches angle 1 as above. This follows by differentiating the divided difference: for \(|u-v|\le c_3(|u|+d_0)\) use the local derivative bounds on \(\mathfrak z'_X\), otherwise divide the bounded slopes by \(|u-v|\). Thus in this exceptional neighborhood its localized \(H^{1/2}\)-norm in \(u\) is at worst polynomial in \(L\) (recall \(1-t\ge 1/s\) if nonzero). Indeed in the difference quotient split translation lengths below and above \(c_3(|u|+d_0)\); the squared seminorm bound costs \(C(1+|\log d_0|)\). At \(t=1\) itself there is no such issue.

\par\smallskip\noindent\textbf{Smoothing the residual kernels.}
For \eqref{eq:M-21} only the mixed error of \(\widetilde H_{ne}\) versus \(H_{ne}\psi\) has to be paired against \(\alpha-\mu\). The displayed smoothing transform gives the Fourier bound \(Ch^{c_2}(1+|k|)^{-1}\widehat\psi\) for the \(S_0\) term. For the decaying residuals just described, convolving in the explicit variable instead of evaluating at its point changes the kernel, as a test in the interior variable, by \(O(\varepsilon^{-C}h^{c_2}+o_L(1))\) in each localized \(H^{1/2}\)-norm, with exponential decay in the distance between unit cells. Indeed the first and mixed derivative bounds suffice by averaging shifts (the probability has small first moment and small exponential tails); when the explicit argument belongs to a curved same-main entry it starts beyond \(L_X\). To meet the exceptional neighborhood at zero it must thus shift by \(\gtrsim L_X\), gaining exponential smallness against the polynomial localized norm there. The smooth real leading tail gives decaying small errors too by even averaging and exponentially decaying second derivatives. For the imaginary leading terms the only change not exponentially localized is from varying \(w_I \upsilon_J(v)\) in the convolution. In \eqref{eq:M-21} this only pairs against \(z\); summing with \(w_I\) and integrating by parts gains \(\tanh'\) against \(D_z\), so the small-error estimate uses just Lipschitz \(\upsilon_J\). Summing tests in spatial cells uses \eqref{eq:M-20}, Sobolev norm bounds from \(B\), and \(\|\psi*\zeta\|_{-1/2}^2\le C(1+|\log h|)\sum n_j^2\) from the previous argument. For clarity, local \(H^{1/2}\) bounds with exponential envelopes combine with loss at most constant in cell \(\ell^2\): insert a unit-scale partition with bounded overlap and use the difference-quotient norm. This proves \eqref{eq:M-21}.

For \eqref{eq:M-22} the singular diagonal real summand \(S_0\) can only decrease on smoothing. Use the same localized residual estimates (only point values needed), now with any same-curved-label explicit points both outside their wall; the \(\log(2\cosh)\) tail again costs small decaying errors. On a same-variable smoothed diagonal the underlying shifts within a label stay regular uniformly after subtracting \(S_0\), giving the stated uniform self bound.


\par\smallskip\noindent\textbf{Summing the localized errors.}
To make the cell summation explicit, if \(a_j\) and \(b_j\) are the
local negative Sobolev norms of the two fields, the error is bounded by
\(C\epsilon_{\rm sm}\sum_{i,j}e^{-c|i-j|}a_i b_j\).
The convolution operator with matrix \(e^{-c|i-j|}\) is bounded on
\(\ell^2\). Cauchy--Schwarz, the magnetic energy norm equivalence, and
the count bound then give
\(C\epsilon_{\rm sm}(B(x)+B(z)+1)\). The omitted diagonal is treated
separately, with one \(C(1+|\log h|)\) cost per explicit particle;
it is never included in the off-diagonal estimate.

\end{proof}

\begin{corollary}[Unshifted coercivity]\label{cor:07_magnetic_analysis-unshifted}
Consequently, for \(z=0\),
\begin{equation}
 \Re\mathcal D\ \ge\ c B(x)-C_{h,\varepsilon}(L+n_e)
 \label{eq:M-23}
\end{equation}
for small suitable \(h,\varepsilon\). 
\end{corollary}
\begin{proof}
The two anomaly logs subtracted from the action in \eqref{eq:M-21} cost arbitrarily small positive fractions of \(B(x)\) plus constants depending on parameters for their real parts (a lower bound suffices). Indeed use the \(H^1\) bounds on main and decaying real upper bounds on far, with \eqref{eq:M-20} and \(x=\Re\alpha-\mu+\psi*\zeta\). The imaginary correction is exactly \(+(\Im\mathfrak t_n,z)\).

Choose \(h\) after \(\varepsilon\) so that the smoothing error is absorbed by the positive real form. Equation~\eqref{eq:M-22} and the diagonal bound leave at most \(C_{h,\varepsilon}n_e\); the bounded anomaly terms are covered by \(C_{h,\varepsilon}L\).
\end{proof}

\subsection{The common interval}

A coarse tilt first bounds the number of far particles in the common
interior. For the remaining configurations, a localized square completion
then extracts the marked-pair gain.

For the remaining estimates assume polynomially bounded \(n_e,|M|,\ell_e\), in fixed powers of \(L\). Our tilts are smooth density derivatives supported strictly within their wall intervals by macroscopic buffers (of size depending on parameters), with their \emph{primitives} compactly supported there, and polynomial fixed-derivative bounds. They depend only on the frozen data, not the variable real interior trial fluctuations. When a primitive used below is constant through the center rather than zero there, we discuss its implementation separately below.


\begin{lemma}[The coarse magnetic tilt]\label{lem:07_magnetic_analysis-coarse}
Take symmetric nonnegative smooth probes \(\varphi_+(u)=\varphi_-(-u)\) of height order 1, rescaled bumps in \(\pm[(1-2\delta)L_X,(1-\delta)L_X]\) for fixed small \(\delta>0\). Let \(I_\varphi=\int\varphi_+=\int\varphi_-\asymp\delta L_X\), and define the nonnegative probe drop by \[ U=I_\varphi^{-1}\int f_R^\psi(u)(\varphi_-(u)-\varphi_+(u))\,du. \] Here \(f_R^\psi\) is the weighted smeared far count to the right and is nonincreasing. A coarse tilt is \(D_z=0,\ z_X=t_0 U(\varphi_+-\varphi_-)\). For small fixed \(t_0>0\) it gives
\begin{equation}
 \Re\mathcal D\ \ge\ c B(x)+c_\delta U^2 L-C_{h,\varepsilon}(L+n_e).
 \label{eq:M-24}
\end{equation}

\end{lemma}
\begin{proof}
Indeed, on the straight comparison contours (main \(X\) at \(\theta_\infty\)) the term \(-(x,\Im H z)\) gives
\begin{equation}
 -2p\int Dz_X-2\pi\int f_R^\psi z_X-(x,J_o z)
 \qquad(D_z=0,\ \int z_X=0).
 \label{eq:M-25}
\end{equation}
The main sign matrix has upper entry \(p\); each elementary far row, less the corresponding main row, has difference of its sign coefficients against the \(X,Y\) columns equal to \(4\beta=\pi\), directly by the far-row transform formulas. The residual \(J_o\) is the straight odd kernel after subtracting its sign tail. The log-factor Fourier formulas give its decaying multiplier bounds, uniformly here since against main columns we're away from real log poles in the straight comparison at fixed parameters (constants polynomial in \(1/\varepsilon\)); near zero frequency these remainders are uniformly bounded also as \(\varepsilon\downarrow0\). For replacing a straight imaginary kernel by the curved kernel against tilts in macroscopic deep regions away from zero, the localized residual after the extra \(\upsilon\)-terms with \(\tanh\) is negligible: where the two arguments are close both are exponentially-asymptotic to straight with derivatives (divide out the self singularity first); elsewhere use the exponential difference-tail estimates with one derivative in each variable. Schur bounds then control the operator against \(H^{-1/2}\)-fields by an exponentially small norm, localizing the tilted variable by a slightly larger deep cutoff. In the extra \(\tanh\) terms \(w_I \upsilon_J(v)\) pairs against weighted-total \(x\) via \(D\), with \(\upsilon_J\) small on tilted support. The term \(-\upsilon_I(u)w_J\) vanishes when \(D_z=0\) (and for a deep-supported \(D_z\) gains exponential smallness against \(\tanh'\)). This justifies \eqref{eq:M-25} up to negligible errors here. The far-count term is exactly \(2\pi t_0 U^2 I_\varphi\). All mixed terms besides the \(f_R^\psi\)-gain cost a fraction of \(B(x)\) plus \(O(t_0^2 U^2\delta L)+o(L)\) by these estimates and low-frequency concentration of the rescaled probes. The cost of \(B(z)\) is similarly bounded; the anomaly term is favorable up to negligible errors by its tails. This proves \eqref{eq:M-24}, taking \(h\) sufficiently small as indicated. In particular for bounded \(U\le U_0\) only \(O(U_0)\) far particles lie in \([-(1-2\delta)L_X,(1-2\delta)L_X]\).

For a particle in the interval between the probe supports,
its contribution to the difference of the normalized averages is bounded
below by a positive constant once \(L\) is large. The exponentially
small smoothing tails do not alter this lower bound. Thus positivity of
the counts gives the asserted \(O(U_0)\) bound on particles there.

\end{proof}

The common and shoulder intervals are shown in Figure~\ref{fig:magnetic-geometry}.
For \(U\le U_0\), localize the sharp tilt into the common region and the two shoulders. Use a smooth square partition of unity in \(u\), including two nonnegative symmetric masks \(w_c,w_s\): the first is 1 in the common bulk (say \(C\delta L_X<|u|<(1-C\delta)L_X\)), supported away from zero and the common walls by \(\delta L_X\), and the second is 1 in the shoulder bulk \(L_X+C\delta L<|u|<L_Y-C\delta L\), supported between these walls away from either by \(\delta L\). Here \(C\) is a fixed sufficiently large constant; the second can be omitted when the interval between walls is sufficiently narrow. Use disjoint separated supports, derivatives \(O_{\delta,j}(L^{-j})\), with remaining masks varying only at distances \(\gtrsim\delta L_X\) from zero too. Such partitions follow by smoothed cosine/sine switches.


\begin{lemma}[Localization of the real energy]\label{lem:07_magnetic_analysis-localization}
For each mask \(w\), set \(D_{y^w}=wD\) and \(y_I^w=wx_I\) for \(I\ne Y\), and define \(y_Y^w=(wD)'-\sum_{I\ne Y}w_Iwx_I\). If \(y^w\) denotes the resulting field, then \[ \sum_w B(y^w)=B(x)+o(1)B(x),\qquad B(y^w)\ge0. \] On the common and shoulder masks, \(B(y^w)\) may be replaced by its straight-contour form \(B_\infty(y^w)\) with error \(o(1)B(x)\). 
\end{lemma}
\begin{proof}
Indeed the real linear-growth tail \(-w_Iw_J|u-v|\) contributes just \(\int D^2\). In the remaining decaying kernels the extra term from differentiating masks in the main \(Y\) charge, of type \(w' D\), is negligible by Sobolev bounds (those decaying kernels are bounded forms on \(H^{-1/2}\); for main \(Y\) entries as needed here this follows already from log diagonal plus smooth localized terms with one derivative in each variable and the harmless \(|u-v|\) cusp). All other errors use the kernel times \(1-\sum w(u)w(v)\), vanishing quadratically at the diagonal, with uniformly integrable bounds \(O_\delta(L^{-1})\) also after up to one derivative in each variable. Indeed this factor kills diagonal log derivative singularities, and where both arguments are near the exceptionally varying central scales the masks are constant (off diagonal at large separation only first graph derivatives are needed). The identity
\[
 1-\sum_w w(u)w(v)=\frac12\sum_w(w(u)-w(v))^2
\]
makes the diagonal cancellation explicit. A differentiated error kernel
has integrable bounds in both variables; applying Schur's test after
half a derivative on each side bounds its pairing on
\(H^{-1/2}\times H^{-1/2}\). The terms involving \(w'D\) are
\(o(1)B(x)\), because \(\|w'\|_\infty=O_\delta(L^{-1})\) and
\(\|D\|_2^2\le CB(x)\). Each localized form is positive by the
magnetic norm equivalence. On \(w_c,w_s\) we may also replace it by the straight \(B_\infty\) with negligible errors using again the log ratios after diagonal division and the decaying residual estimates (real growth tail unchanged). Write the localized real fields on these masks as \(y^c,y^s\).


\end{proof}

\begin{lemma}[The common-region trial gain]\label{lem:07_magnetic_analysis-common}
For the unscaled common-region tilt, set \(z_X=A_0\operatorname{sign}(u)w_c^2,\ z_Y=-z_X,\ A_0=6\beta/(2\kappa)\), denoting these components by \(z^c\). The reserve factor \(v_0\) is introduced in Proposition~\ref{prop:07_magnetic_analysis-sharp}. Then
\begin{equation}
 B_\infty(y^c)-B(z^c)-(x,\Im H z^c)+(\Im\mathfrak t_n,z^c)
 \ \ge\ \int w_c^2\,du-o(B(x)+L)
 \label{eq:M-26}
\end{equation}
with errors understood as terms having coefficient tending to zero. 
\end{lemma}
\begin{proof}
In \eqref{eq:M-25} the far part is favorable by monotonicity, and the anomaly gives \(6\beta A_0\int w_c^2+o(L)\). The odd remainder costs \(o(\sqrt L)\sqrt{B(x)}\) by Fourier bounds and low-frequency concentration. The cost \(B(z^c)\) is \((\kappa-p^2)A_0^2\int w_c^4+o(L)\) by the straight main symbol. At small frequency the real \(B_\infty\) symbol in variables \(D,x_{I\ne Y}\) has the \(D\) coefficient 1 and mixed \(D\) terms vanishing at zero; the complementary block stays strictly positive by the real norm bound. To spell out the square completion, let \(\mathcal B(k)\) be the
real energy symbol in these variables, with the convention that its
\(D,D\) entry tends to one. For
\(f_L(u)=\operatorname{sign}(u)w_c(u)\), minimizing the quadratic
form against the linear test at each frequency gives
\[
 B_\infty(y^c)-2pA_0\int (w_cD)f_L
 \ge -\frac{p^2A_0^2}{2\pi}
     \int |\widehat f_L(k)|^2
                [\mathcal B(k)^{-1}]_{DD}\,dk.
\]
The inverse entry tends to one at zero, since the mixed block vanishes
there and the complementary block is strictly positive. Its growth
away from zero is polynomial by the energy norm equivalence. The
smooth scaled function \(f_L\) has Fourier mass concentrated at
frequencies of order \(L^{-1}\), with arbitrarily rapid decay outside
that scale. Splitting the integral at a small fixed frequency and then
letting that cutoff tend to zero shows that the right side equals
\(-p^2A_0^2\int w_c^2-o(L)\). Finally \(0\le w_c\le1\) and \(\kappa>p^2\), so the sum of the two negative quadratic costs is at least \(-\kappa A_0^2\int w_c^2-o(L)\). Since \(A_0=3\beta/\kappa\), \(6\beta A_0-\kappa A_0^2=9\beta^2/\kappa=1\), proving \eqref{eq:M-26}.


\end{proof}

\subsection{The shoulders}

Only the $Y$ main particles remain in the shoulders. Eliminating their
real and imaginary responses leaves a reduced kernel on every other
label. Its positivity and self cost will control the shoulder activities.

\begin{lemma}[The shoulder trial gain]\label{lem:07_magnetic_analysis-shoulder}
In the shoulder calculation let \(e\) index \emph{all labels except main \(Y\)}, and use straight translation-invariant kernels below. Set, as matrix kernels,
\[
 \mathcal H=\Re(H_{ee}-P),\qquad
 \widehat P=\widehat H_{eY}\widehat H_{YY}^{-1}\widehat H_{Ye}
 \quad(k\ne0)
\]
with \(\Re P\) chosen to have real tail \(-w_Iw_J|u-v|\) plus decaying terms. Put \(K_1=(\Im H_{eY})\partial,\ H_p=-H_{YY}\partial^2\), where \(H_p\) denotes the positive scalar primitive symbol \(k^2\widehat H_{YY}\) (including its positive limit at zero). Take
\[
 a=w_s(\psi*\zeta)_e,\qquad
 V_0=-H_p^{-1}K_1^{\mathsf t} a,\qquad
 D_{z^s}=w_s V_0,\quad z^s_X=0,\quad z^s_Y=(w_s V_0)' .
\]
The multipliers to \(V_0\) are exponentially decreasing analytic symbols on some strip around real frequency, by the far-row and straight main formulas (\(X\) heights at least \(1/2\), far interactions against main \(Y\) regular with margin). All needed derivatives of \(V_0\) are thus polynomially bounded under count truncations. Moreover \(a-w_s x_e\) is exponentially small in the negative norm (only the tail mean of main \(X\)), so these fields have the required Sobolev sizes bounded by \(C\sqrt{B(x)+1}\), with \(C\) uniform at small \(\varepsilon\) and sufficiently large \(L\). Then
\begin{equation}
 B_\infty(y^s)-B(z^s)-(x,\Im H z^s)+(\Im\mathfrak t_n,z^s)
 \ \ge\ \tfrac12(a,\mathcal H a)-o(B(x)+1).
 \label{eq:M-27}
\end{equation}

\end{lemma}
\begin{proof}
Indeed the anomaly pairs against a derivative supported deep in each tail with vanishing integral there. The curved imaginary corrections are negligible as explained after \eqref{eq:M-25}; on straight kernels the mixed term is \(-(w_s x_e,K_1 V_0)\) up to negligible commutators (the differentiated kernels are smooth exponentially decreasing with derivatives). Also \(2 B(z^s)=(w_s V_0,H_p w_s V_0)\le(V_0,H_p V_0)+o(B(x)+1)\) by the same square partition argument on main \(Y\) alone. This recovers plus half the square of the imaginary coupling with \(H_p^{-1}\). In \(B_\infty(y^s)\) minimizing over the remaining real total primitive at each frequency gives the real Schur subtraction \((\Re H_{eY})H_{YY}^{-1}(\Re H_{Ye})\) from \(\Re H_{ee}\); here both are combined before integration, via the neutral form symbols, so no zero-frequency subtraction constants enter. The resulting bounded symbol (order \(-1\) at infinity) acts on \(w_s x_e\), replaceable by \(a\) as above. Combining gives \eqref{eq:M-27}. The cross term in \(B(z^c+z^s)\) is negligible by support separation and the zero common total, using exponentially decaying real residual kernels.


The sign in this completion is essential: if \(b=K_1^{\mathsf t}a\),
then \(V_0=-H_p^{-1}b\) gives
\[
 -\tfrac12(V_0,H_pV_0)-(b,V_0)
 =\tfrac12(b,H_p^{-1}b).
\]
Adding this positive term to the real Schur complement produces
\(\operatorname{Re}(H_{ee}-H_{eY}H_{YY}^{-1}H_{Ye})\), rather than
the Schur complement of \(\operatorname{Re}H\) alone.

\end{proof}

\begin{lemma}[Positivity and self cost of the reduced shoulder kernel]\label{lem:07_magnetic_analysis-positive}
For the straight asymptotic labels and all sufficiently small
\(\varepsilon>0\), the tail-normalized kernel \(P\) above satisfies
\begin{equation}
 \mathcal H_{IJ}(r)\ge-C\varepsilon e^{-c_2|r|},\qquad
 \operatorname{Re}P_{II}(0)\le-w_I^2\pi/2.
 \label{eq:M-28}
\end{equation}
The first inequality is entrywise off point singularities. Its constants
are uniform as \(\varepsilon\downarrow0\).
\end{lemma}
\begin{proof}
\par\smallskip\noindent\textbf{The reduced kernel.}
We first identify the reduced kernels, and then estimate their off-diagonal and diagonal parts separately. Use \(x=\beta k\). Define
\[
 F_y(r)=2\operatorname{atanh}(\sin(\pi y/6)\operatorname{sech}(\pi r/(6\beta))),
 \qquad G_y=F_{y+1}-F_y+F_{y-1}.
\]
For \(|y|\le3\), \(F_y\) has symbol \(2\pi\sinh(yx)/(k\cosh(3x))\): differentiate in \(y\) and use the elementary sech transform by residues (imaginary shifts inside the first poles); extend to the endpoints by log integrability. Beyond that range use the foldings indicated by the sine.

For two elementary constituents let \(d\) be their absolute height difference in \(\beta\) units. Then the following give their contributions to \(\mathcal H\):
\begin{itemize}
\item across groups, \(G_{d-4}\);
\item within a group, \((G_{2+d}+G_{2-d})/2\), plus only in group \(X\) a correction with symbol
\[
 \frac{2\pi}{k}\frac{c_*(k)\sinh x\cosh(dx)}{4\cosh x\cosh(2x)\cosh(3x)} .
\]
\end{itemize}
Indeed divide all \(H_{IY}\) symbols by \(2\pi c_*/(k\sinh(4x))\); the contributions are \(e^{(3-a)x}\) for \(X\) elementary at height \(a\), \(e^{(-4+j)x}\cosh x\) for \(Y\) far at height \(-j\), with main \(Y\)-entry \(\cosh(3x)\). For a cross pair, before folding at \(d+h=8,\ h=0,1,2\), its real \(H\)-numerator in the same normalization is \(\cosh((3-d)x)\). The subtraction uses \(\cosh((7-d)x)\cosh x/\cosh(3x)\), giving just the unfurled \(G_{d-4}\) expression by product-to-sum. Each crossing replaces a \(\cosh((4-d-h)x)/\sinh(4x)\) by its reflected expression, adding \(2\sinh((8-d-h)x)\) with the same coefficient; this is precisely the fold correction of \(F_{d+h-5}\). For a same-group pair before crossings at \(d=1,2\), the real symbol divided now by \(2\pi/k\) is \(c_* \cosh(3x)\cosh(dx)/\sinh(4x)-\sinh(dx)\); the last term corresponds to folding \(F_{3+d}/2\). Subtracting \(c_* \cosh^2(x)\cosh(dx)/(\sinh(4x)\cosh(3x))\) gives the stated \(Y\) expression, with crossings again exactly folds (additional \(+\sinh((d-1)x),-\sinh((d-2)x)\) respectively). For \(X\) remove the \(\cosh^2(x)\) in that subtraction, giving the correction. This covers the whole indicated ranges. Real tails cancel at infinity, excluding affine ambiguities; the prescription of \(\Re P\) with two derivatives given by the symbols is smooth analytic with exponentially decaying second derivatives as used above (heights here within fixed decay margins).

\par\smallskip\noindent\textbf{Positivity at zero far-contour parameter.}
All resulting \emph{fused} entries are nonnegative at \(\varepsilon=0\), off point singularities. For the extra correction at $|d|\le3$, use the three factors
\[
 \frac{\sinh x}{k\cosh x},\qquad
 \frac{c_*}{\cosh(2x)},\qquad
 \frac{\cosh(dx)}{\cosh(3x)}.
\]
Each has a positive inverse transform as a measure: use averaging and real parts of imaginary-shifted sech (for \(c_*=2\cosh x-1\), the shifted real sech dominates \(\cos(\pi/4)\) times the unshifted). When one partner is a double, pair differences \(d,d+1\) signed; combining the two uses \(c_*\cosh(x/2)=\cosh(3x/2)\), also positive by the same test for \(|d+1/2|\le3\), which holds at zero parameter. Any perturbative deficit here is \(O(\varepsilon)e^{-c_2|r|}\) by the individual analytic symbols with decay margins.

For the \(G\)-terms note \(G_y\ge0\) for \(0\le y\le 21/8\): with \(t_1=\pi y/6,\ q_0=\operatorname{sech}(\pi r/(6\beta))\), the hyperbolic tangent addition comparison reduces to
\(\sqrt3\ge1+q_0^2(\sin^2 t_1-1/4)\), true since \(\sin^2(7\pi/16)+3/4<\sqrt3\). This handles the singleton cross pairs, while double cross partners combine by \(G_v+G_{v+1}=F_{v-1}+F_{v+2}\ge0\) at \(-1/2\le v\le3\). For like-group singleton pairs only \(d=0\) or \(d\in[17/8,21/8]\) are needed. In the latter case \(G_{2+d}=G_{2-y}\ge G_y=-G_{2-d}\), \(y=d-2\): compare odd Taylor powers of atanh at \(t_1=\pi y/6\in[\pi/48,5\pi/48]\) and \(\pi/3-t_1\). The linear difference favors the latter; for third powers use \(\sin^3\alpha=(3\sin\alpha-\sin(3\alpha))/4\), whose triple-angle contribution is identical. For odd orders \(m\ge5\), the difference (apart from common coefficients) is
\[
 \cos^m t_1-\sin^m(\pi/3-t_1)+2\sin^m(\pi/6-t_1)-\sin^m(\pi/6+t_1)+\sin^m t_1
 \ \ge\ .945^m-.833^m-.754^m>0.
\]
The last strict inequality holds at \(m=5\), where the difference exceeds \(0.1088\). Dividing by \(.945^m\) shows that it persists for every larger \(m\), since both ratios \(.833/.945\) and \(.754/.945\) are less than one. For a same-group double combine at signed differences \(d,d+1\); the four \(G\)'s become \(F_{1+d}+F_{4+d}+F_{-d}+F_{3-d}\). The necessary lower differences (double heights counted upwards in \(X\), or in absolute value in \(Y\)) at zero parameter are \(d\in[2,5/2]\) or \(d\in[-1,0]\). In the first case \(F_{1+d}\ge F_d,\ F_{3-d}\ge F_{d-2}=-F_{4+d}\); in the second all terms are positive or zero. 

\par\smallskip\noindent\textbf{Moving the far double.}
For small \(\varepsilon>0\) only combinations involving the moved \(X\)-double can change. Cross partners there continue positive: for \(1\le v\le3\)
use the preceding identity, while for \(3\le v\le4\),
\(\sin(\pi(v-1)/6)\ge-\sin(\pi(v+2)/6)\), so that
\(F_{v-1}+F_{v+2}\ge0\). Thus the whole range \(1\le v\le4\)
is covered; in the just indicated same-group comparisons, every loss costs only \(C\varepsilon e^{-c_2|r|}\) by differentiation wherever the comparison may fail (e.g. \(d\) near \(5/2\)), away from negative log divergences. Explicitly the case \(d\in[2+\varepsilon,5/2+\varepsilon]\) obeys the same inequalities until \(5/2\), and the other moving lower difference is \(d=-1/8+\varepsilon\). This proves the first bound in \eqref{eq:M-28}.

\par\smallskip\noindent\textbf{The diagonal self cost.}
Finally by tail normalization
\[
 \Re P_{II}(0)=w_I^2\int_0^\infty
 \left[\frac{2 c_*(x/\beta) A_I(x)^2}{x\sinh(4x)\cosh(3x)}-\frac1{2x^2}\right]dx
\]
where \(A_I=1\) or \(\cosh x\) for \(X\) or \(Y\) respectively, times \(\cosh(x/2)\) for a double. We have
\[
 4x(2\cosh x-1)\cosh^2 x\cosh^2(x/2)(1+4x^2)
 \ \le\ \sinh(4x)\cosh(3x),\qquad x\ge0.
\]
Indeed the coefficient condition by product-to-sum is
\((2n+1)(b_n+4(2n)(2n-1)b_{n-1})\le7\cdot49^n+1\),
\(b_n=\delta_{n0}+3+2\cdot4^n+9^n+16^n\), second summand absent at zero. For \(n=0,1,2,3\) the left sides are \(8,300,10500,347172\); thereafter both left summands grow by a factor at most \(16(2n+3)(2n+2)/(2n(2n-1))\) from \(n\) to \(n+1,\ n\ge3\). This factor is at most \(38.4\) for \(n\ge3\), whereas \((7\cdot49^{n+1}+1)/(7\cdot49^n+1)>48\); the four checked cases therefore initiate the induction. Thus the bracket is bounded above by \(-2/(1+4x^2)\), and \(\int_0^\infty 2/(1+4x^2)\,dx=\pi/2\), proving the second inequality in Equation~\eqref{eq:M-28}.


\end{proof}

\subsection{Combining the trial fields and summing explicit particles}

We combine the common and shoulder fields while reserving enough positive
energy to absorb smoothing errors. The explicit particles will then be
summed with their actual activities and factorial measures.

\begin{proposition}[The sharp magnetic trial estimate]\label{prop:07_magnetic_analysis-sharp}
Suppose \(U\le U_0\). Fix sufficiently small \(\varepsilon>0\), choose \(h=\varepsilon^{K_2}\) with \(K_2\) sufficiently large, and put \(v_0=1-\sqrt\varepsilon\). For the tilt \(v_0(z^c+z^s)\),
\begin{equation}
\begin{split}
 \Re\mathcal D \ \ge\ &
 c_{\varepsilon,h}\{B(x)+\ell_e+(M^2+n_e^2)/L\}
 +2L_X-o_\varepsilon(1)L-O(\delta L)-o(L)\\
 &-\sum_{v\in\zeta}
 \left[\frac{v_0}2 w_s(v)^2\Re P_{I(v)I(v)}(0)
      +C(1+|\log h|)(1-v_0 w_s(v)^2)\right].
\end{split}
 \label{eq:M-29}
\end{equation}
Use zero for an unnecessary \(P\) entry in this notation. Constants depending on \(h,\varepsilon,\delta\) have been absorbed in the little-oh in \(L\). Here \(c_{\varepsilon,h}>0\) may be chosen independently of small fixed \(\delta\), and \(o_\varepsilon(1)\to0\) in the indicated parameter order.


\end{proposition}
\begin{proof}
To use \eqref{eq:M-26},\eqref{eq:M-27} simultaneously while retaining control, fix small far parameter \(\varepsilon>0\), take \(h=\varepsilon^{K_2}\) for a sufficiently large fixed \(K_2\), set \(\eta_0=\sqrt{\varepsilon}\), and \emph{scale both tilts} by \(v_0=1-\eta_0\). Keep \(v_0 B(x)\) for localization; the tilt cost is no worse than scaling the old costs by \(v_0\), leaving a real-energy reserve. The smoothing errors of \eqref{eq:M-21} can use small portions of that reserve (and \(o_\varepsilon(1)L\)), since the unscaled \(B(z)\le C(B(x)+L)\). In \(\Delta_H\) split off the piece with entry weights \(v_0 w_s(u)w_s(v)\) on the explicit positions (same convention for missing diagonals). On the complement \eqref{eq:M-22} costs \(o(\eta_0)B(x)+o(1)\), plus the diagonal omission bound
\[
 -C(1+|\log h|)\sum_{v\in\zeta}(1-v_0 w_s(v)^2).
\]
For the piece split off only \(e\) labels intervene. All curvature changes there are negligible even after smoothing, since straight approximation holds exponentially well on deep support (with log ratios at the main diagonal), with buffers to any exceptional smearing regions where only pointwise residual bounds are needed. After combining with \(v_0(a,\mathcal H a)/2\) the result is
\[
 \frac{v_0}2(w_s\zeta_e,\mathcal H w_s\zeta_e)_{\rm off}
 -\frac{v_0}2\sum_{v\in\zeta_e}w_s(v)^2\Re P_{I(v)I(v)}(0)
\]
up to negligible errors or \(o(\eta_0)(B(x)+1)\). Indeed one can replace \(a\) by \(\psi*(w_s\zeta_e)\) in that quadratic by slow variation, with error at most \(C_{\varepsilon,h,\delta} L^{-1}\sum n_j^2\) using the decaying kernel \(\mathcal H\) (at most local log singularities, as proved in Lemma~\ref{lem:07_magnetic_analysis-positive}). The remaining averaging errors use \(\Re(P-\psi P\psi)\), \(O(h^{c_2})\) with exponential spatial decay, since after two derivatives the symbol of \(\Re P\) is analytic and exponentially decreasing, uniformly here.

Thus \eqref{eq:M-20} absorbs the first lower-bound error by the reserve. 
In this argument the off-diagonal deficit is at most
\(C\varepsilon\sum_j n_j^2\). By Equation~\eqref{eq:M-20}, its
energy coefficient is \(O(\varepsilon(1+|\log h|))=o(\sqrt\varepsilon)\).
The smoothing errors are smaller still when \(K_2\) is large.
These estimates leave a fixed positive fraction of the
\(\sqrt\varepsilon B(x)\) reserve. Applying Equation~\eqref{eq:M-20}
to a further fraction of that reserve gives the tail and count terms in
\eqref{eq:M-29}. Finally \(\int w_c^2=2L_X-O(\delta L)\), and scaling the gain
by \(v_0\) loses only \(O(\sqrt\varepsilon L)\).

\end{proof}

\begin{corollary}[Summation of explicit particles]\label{cor:07_magnetic_analysis-activities}
After applying Proposition~\ref{prop:07_magnetic_analysis-sharp}, the sum and integrals over explicit particles with \(U\le U_0\) cost at most
\begin{equation}
 \exp\big((.44+o_\varepsilon(1))\,2(L_Y-L_X)
       +C_{\varepsilon,h}\delta L+o(L)\big)
 \label{eq:M-30}
\end{equation}
besides the negative bulk cost and any reference fluctuation cost. 
\end{corollary}
\begin{proof}
Indeed in deep shoulders there are three singleton types and two double types (no main \(Y\)); \((\mu,R\delta_v)=o(1)\), curve Jacobian modulus \(1+o(1)\) where applicable, and the deep extra log penalty \(C(1+|\log h|)\eta_0=o(1)\) per particle. By \eqref{eq:M-28} the limiting sum of grand activities per unit length there is bounded by
\(3(d/c)e^{-\pi/4}/(2\pi)+2\cdot4c^2(d/c)^2 e^{-\pi}/(2\pi)=0.429757\ldots<0.44\).
All boundary and transition buffers have only \(O(\delta L)\) volume; two infinite ends also use the \(\ell_e\) damping. Common interior far counts are bounded by \eqref{eq:M-24}; thus factorial integration with the indicated per-particle constants gives \eqref{eq:M-30}. The polynomially many allowed \(M\)'s cost only \(e^{o(L)}\). In using coarse bounds, summing with \(e^{C_{\varepsilon,h} n_e}\) and part of the retained damping costs at most \(e^{O_{\varepsilon,h}(L)}\); \(U_0\) can be chosen sufficiently large afterwards.


For example, for a species of activity \(a\) on a set of length \(V\),
the factorial sum is
\[\sum_{m\ge0}(aV)^m/m!=e^{aV}.\] This identity applied species by
species explains both the shoulder contribution and the buffer-volume
bound. On an infinite end, replace \(V\) by
\(\int_0^\infty e^{-c_{\varepsilon,h}r}\,dr\), retaining the
\(\ell_e\) penalty. In the common interior the bounded number of far
particles contributes at most a fixed power of \(L\), hence \(e^{o(L)}\).

\end{proof}

\subsection{Selecting marks for the polygon moment}

We also use the planar polygon tail
\begin{equation}\label{eq:planar-tail}
 \sum_{[P]:\operatorname{diam}P>R}\rho_{\rm v}^{|P|}
       \le C(1+R)^{-2},\qquad R\ge0,
\end{equation}
from \cite[Theorem~10.5]{Electric}. Its classes are unrooted,
unoriented polygons modulo triangular-lattice translations, with
vertex length. This positive estimate is used only to dispose of
polygons too short for the geometric marked-pair selection. It is
distinct from the reference-law and contour inputs used to estimate the marked integral.

\begin{lemma}[Reduction to marks with macroscopic margins]\label{lem:07_magnetic_analysis-marks}
For the polygon second-moment estimate, it suffices to implement the trial bounds under the additional margin \(L_X\ge c_1 L\) and, when \(t\ne1\), \(s(1-t)\ge e^{c_4 L}\) for a fixed sufficiently small \(c_4>0\). This latter margin matters for tilting primitives through the central region. More precisely, a uniform marked-pair estimate
\[ \mathcal C(p,p')\le D^{o(1)}s^{-4/3}(D/s)^{2/3} \]
for the resulting configurations with \(b\asymp D\) implies the required
second-length moment bound at dyadic diameter \(D\).
\end{lemma}
\begin{proof}
Fix a dyadic diameter bin $D/2<\operatorname{diam}P\le D$. We
retain a fixed positive fraction of the ordered vertex pairs on every
sufficiently long polygon in this bin, then compare their weights with
marked cylinder correlations.

\par\smallskip\noindent\textbf{Short polygons and retained pairs.}
By Equation~\eqref{eq:planar-tail}, polygons of length at most
$D^{4/3}$ contribute at most
$CD^{-2}D^{8/3}=O(D^{2/3})$ to the second-length mass in this bin.
On a longer polygon, retain the ordered pairs whose separation exceeds
$D^{1/2}$ and whose displacement has perpendicular component greater
than $D^{1/8}$ in absolute value to each of the three triangular
lattice directions. For each first vertex, the excluded disk contains
$O(D)$ lattice vertices, and the three excluded strips within distance
$D$ contain $O(D^{1+1/8})$. Since $|P|>D^{4/3}$, at least half of
its ordered vertex pairs remain once $D$ is large.

\par\smallskip\noindent\textbf{A row determined by the two ports.}
Choose a triangular angular basis whose positive $60^\circ$ cone
contains the directed displacement. At each endpoint triangle, at
least one of the two occupied sides has an allowed direction, since
only one of the three side directions is excluded. Move the endpoint
to such an occupied port, using a fixed rule if there is a choice.
These moves have bounded length and preserve the retained margins.
Write $e_1,e_2$ for the chosen unit basis sides and $h_p,h_{p'}$ for
the positive half-side vectors at the resulting ports. Then
\[
 p'-p-h_p-h_{p'}=m_1e_1+m_2e_2,\qquad
 m_1,m_2\in\mathbb Z_{\ge0},\quad m_j\gg D^{1/8}.
\]
The left side joins vertices of the triangular tiling, which proves
integrality. The angular margin gives the two positive lower bounds.
Travel from $p$ along $h_p$, then along the $m_1+m_2$ complete sides
in a fixed order, and finally along $h_{p'}$ to $p'$. This gives a
row with $s=m_1+m_2$ ordinary bands and
$s\asymp|p-p'|$.

Choose among the finitely many rotations and reflections so that the
larger coefficient is assigned to the upper site height. Let
$e_{\rm low}$ be the other basis direction. Continue from $p'$ along
its remaining half-side, then along $b$ complete sides in direction
$e_{\rm low}$, and finish with a half-side parallel to $h_p$. The
period vector is
\[
 T=m_1e_1+m_2e_2+b e_{\rm low}+2h_p+2h_{p'}.
\]
All increments lie in the same positive cone, so the repeated row is
simple and $|T|\ge T\cdot e_{\rm low}\ge b$. Choose $b$ to be the
smallest integer at least $KD$ for which $s+b+2$ is even, with $K$
a fixed sufficiently large constant. Then $b\asymp D$, $s\le b$,
and $|T|>2D$. No two points of $P$ can differ by a nonzero integer
multiple of $T$, so its cylinder projection is injective. This
choice of row and period depends only on the port data, the chosen
basis, and $D$. Figure~\ref{fig:marked-row} shows the endpoint
half-sides, a fixed order of the ordinary bands, and the complementary
part of the period.
\begin{figure}[htbp]
\centering
\begin{tikzpicture}[x=0.98cm,y=0.98cm,>=Latex,font=\small,
  ordinary/.style={circle,draw,fill=white,inner sep=1.7pt},
  marked/.style={circle,fill=black,inner sep=2.1pt},
  vertex/.style={circle,fill=black!35,inner sep=1pt}]
 \coordinate (p) at (0,0);
 \coordinate (a) at (0.5,0);
 \coordinate (bend) at (2.5,0);
 \coordinate (c) at (4,{1.5*sqrt(3)});
 \coordinate (pp) at (4.25,{1.75*sqrt(3)});
 \coordinate (d) at (4.5,{2*sqrt(3)});
 \coordinate (e) at (9.5,{2*sqrt(3)});
 \coordinate (pt) at (10,{2*sqrt(3)});

 \draw[thick,blue!65!black] (a)--(bend)--(c);
 \foreach \j in {0,1} {
   \node[ordinary,draw=blue!65!black] at ({1+\j},0) {};
 }
 \foreach \j in {0,1,2} {
   \node[ordinary,draw=blue!65!black]
     at ({2.75+0.5*\j},{(0.25+0.5*\j)*sqrt(3)}) {};
 }
 \node[blue!65!black,above] at (1.5,0.13) {$X$};
 \node[below] at (1.5,-0.08) {$m_1e_1$};
 \node[blue!65!black,rotate=60,above] at (3.15,1.3) {$m_2e_2$};

 \draw[black!25] (-0.5,0)--(p);
 \draw[thick,purple!70!black,->] (p)--(a);
 \draw[thick,purple!70!black,->] (c)--(pp);
 \draw[thick,purple!70!black] (pp)--(d);
 \draw[thick,purple!70!black] (e)--(pt);
 \draw[black!25] (pt)--(10.5,{2*sqrt(3)});
 \node[purple!70!black,above] at (0.23,0.12) {$h_p$};
 \node[purple!70!black,left] at (4.08,{1.68*sqrt(3)}) {$h_{p'}$};

 \draw[thick,orange!75!black] (d)--(e);
 \foreach \j in {0,1,2,3,4} {
   \node[ordinary,draw=orange!75!black]
     at ({5+\j},{2*sqrt(3)}) {};
 }
 \node[orange!75!black,above] at (7,{2*sqrt(3)+0.12})
   {$Y$: $b$ bands in direction $e_1$};

 \foreach \pos in {a,bend,c,d,e} \node[vertex] at (\pos) {};
 \node[vertex] at (1.5,0) {};
 \foreach \j in {1,2} {
   \node[vertex] at ({2.5+0.5*\j},{0.5*sqrt(3)*\j}) {};
 }
 \foreach \j in {1,2,3,4} {
   \node[vertex] at ({4.5+\j},{2*sqrt(3)}) {};
 }
 \node[marked,label={left:$p$}] at (p) {};
 \node[marked,label={above left:$p'$}] at (pp) {};
 \node[marked,label={below right:$p+T$}] at (pt) {};

 \draw[densely dashed,->,black!75]
   (0.3,-1.25)--(10.3,{2*sqrt(3)-1.25})
   node[midway,sloped,below,fill=white,inner sep=2pt]
     {$T=m_1e_1+m_2e_2+be_1+2h_p+2h_{p'}$};

 \begin{scope}[shift={(0.2,2.05)}]
   \draw[->] (0,0)--(1,0) node[right] {$e_1$ (low)};
   \draw[->] (0,0)--(0.5,{sqrt(3)/2})
     node[above] {$e_2$ (upper)};
 \end{scope}
\end{tikzpicture}
\caption{Constructing the cyclic cut from an ordered pair of visited
ports. Open circles are ordinary ports at the midpoints of complete
triangular sides; black circles are the marks. The $X$ part joins
$p+h_p$ to $p'-h_{p'}$ using $m_1$ sides in direction $e_1$ and
$m_2$ in direction $e_2$, in a fixed order. The complementary part has
$b$ ordinary sides in the low direction. The displayed mark types have
$h_p=e_1/2$ and $h_{p'}=e_2/2$; either half-side can have either
allowed direction. Identifying $p+T$ with $p$ closes the row on the
cylinder. Taking $|T|>2D$ makes projection injective on each planar
polygon of diameter at most $D$. Counts and lengths are illustrative.}
\label{fig:marked-row}
\end{figure}


The retained pairs have $s\gg D^{1/2}$, and the minority count is
$s(1-t)=\min(m_1,m_2)\gg D^{1/8}$. Consequently
\[
 L_X\ge\frac{\log D}{2\gamma}-O(1),\qquad
 L_Y=\frac{\log D}{\gamma}+O(1).
\]
Every fixed $0<c_1<1/2$ and $0<c_4<\gamma/8$ therefore gives
$L_X\ge c_1L_Y$ and $s(1-t)\ge e^{c_4L_Y}$ for all sufficiently
large $D$. These configurations meet the stated margins.

\par\smallskip\noindent\textbf{Counting the marked occurrences.}
Moving a vertex to an incident port has at most two inverse choices
at each endpoint. The choices of basis, rotation, reflection, and
port type also have bounded multiplicity. Normalize the first port
modulo triangular-lattice translations. Once its lift is fixed, the
projected polygon has a unique planar lift through that port, with
its original vertex weight. Thus the retained pairs account for at
least a fixed multiple of $|P|^2$ marked occurrences, and any one
marked cylinder polygon arises from at most a fixed number of these
occurrences.

For a dyadic displacement scale $S$, let $\mathscr R_{D,S}$ be the
row setups just constructed with $S\le s<2S$, and write
$\mathcal C_{\mathfrak r}$ for the corresponding marked correlation.
There are $O(S^2)$ setups: after fixing the first port type, there
are only $O(S^2)$ possible second-port displacements, and the
remaining choices are bounded. Positivity and the preceding
multiplicity bounds give
\[
 \sum_{\substack{[P]:\ D/2<\operatorname{diam}P\le D\\
                         |P|>D^{4/3}}}
       \rho_{\rm v}^{|P|}|P|^2
 \le C\sum_{\substack{S\ \text{dyadic}\\D^{1/2}\ll S\ll D}}
          \ \sum_{\mathfrak r\in\mathscr R_{D,S}}
                    \mathcal C_{\mathfrak r}.
\]
A bound of the stated form makes each displacement scale contribute
at most $D^{2/3+o(1)}$. The $O(\log D)$ scales and the
short-polygon estimate prove the claimed reduction.
\end{proof}

\section{Magnetic fluctuations and the rare central scale}
\label{sec:07_magnetic_analysis-fluctuations}

The trial bounds have reduced the problem to a reference integral with
fixed integer dimensions. Its Hessian sees the entire wall intervals
on the main contours, including the rare central scale. This differs from the preceding
smoothing estimates, where a same-main explicit particle can reach that
scale only through an exponentially small smoothing tail. We first
control the Hessian and its determinant at the rare scale. We then
choose a coordinate in which the trial fields can be implemented by
small particle motions, adapting the reference and contour arguments
of \cite{Electric} to the unequal magnetic walls.

\subsection{Reference coordinates and the relative Hessian}

We now implement \eqref{eq:M-23},\eqref{eq:M-24},\eqref{eq:M-29} under the parameter margins just stated. Retain the sizes $s_g$ and unit-mass mean densities $\bar\nu_g$ introduced in the trial setup. In this part no interpolation of dimensions, nor continuation in \(M\), is needed: use \(N_g=m'_g\) particles in the reference integral of group \(g\) (only nonnegative integer dimensions occur). All inner trial changes have zero mass separately by group. After the
polynomial count truncation justified in the proof of
Proposition~\ref{prop:07_magnetic_analysis-transport},
\(N_g=s_g+O(L^K)\) for a fixed \(K\), whereas
\(s_g=e^{\gamma L_g}\); hence \(N_g\asymp s_g\).
Phase factors depending only on counts have modulus one.

We may work on subsequences on which \(L_X/L\) and the mixture fraction \(t\) have limits. Call the case \(\lim t<1\), or \(t=1\) identically, regular. In the remaining case, passing to another subsequence, \(t<1,\ t\to1\); call this the rare case and use the scale \(d_0=\beta(1-\theta(0))\) from the curve estimates. Thus \(\log(1/d_0)=O(L)\), and \(s d_0\) grows at least as \(e^{cL}\) for some \(c>0\).

We first specify the cap maps used in the reference coordinates.
For an interval $[-\Lambda,\Lambda]$ choose a fixed $a_0=O(1)$
and an odd map $W:[-\Lambda-a_0,\Lambda+a_0]\to[-\Lambda,\Lambda]$.
At the positive endpoint write
$W(\Lambda+a_0-r)=\Lambda-\omega(r)$, where $\omega$ is quadratic
near zero, increasing with $\omega'(r)\asymp\min(r,1)$, and
$\omega(r)=r-a_0$ for $r\ge1$. Integrating a smooth cutoff between
$r$ and one constructs $\omega$ and fixes $a_0$. Thus $W(v)=v$
away from bounded endpoint neighborhoods and reflects evenly at the
endpoints.

Use first an \emph{unstretched} parameter $v$ per wall interval. In the regular case set the intermediate coordinate \(\xi=u\) on each, and use the cap map \(\xi=W_g(v)\) onto \([-L_g,L_g]\) using the cap construction just given for each separate length. Take the reference probability with density in \(v\)
\[
 \lambda_g^0(v)=C_g\,\eta_g(W_g(v)),\qquad \eta_g=\bar\nu_g .
\]
In the rare case do this on \(Y\) as usual, but on \(X\) set \(\xi=b_X(u)\) by the mean-matching map, use \(W_X\) onto \([-b_X(L_X),b_X(L_X)]\), and take \(\eta_X=\rho\). Write \(A_g\) for the unstretched positive endpoints and \(U_g(v)\) for the resulting map to \(u\). In all cases \(A_g=L_g+O(1)\), \(C_g=1+O(e^{-\gamma L_g})\). The push of the unnormalized density \(\eta_g(W_g(v))\,dv\) to \(u\) agrees with \(\bar\nu_g\,du\) away from or before the bounded-width cap regions (i.e. throughout the interior bulk); the cap and omitted tail masses are \(O(e^{-\gamma L_g})\) with exponential decay outside. Denote the normalized push by \(p_g\), and the reference angular coordinate by \(t_g(v)=\pi\int_{-A_g}^v\lambda_g^0\). The curve estimates show that the unstretched densities have uniform fixed-order relative derivative bounds, sizes comparable to $e^{-\gamma|v|}$, and uniform local analytic bounds off the caps. They extend evenly at the endpoints.

The coordinate maps are thus
\[
 v\ \xmapsto{\ W_g\ }\ \xi\ \longmapsto\ u=U_g(v),
 \qquad
 \xi=u\ \text{except on rare }X,\quad \xi=b_X(u)\ \text{there}.
\]
The contour point is $\mathfrak z_g(u)$, while $t_g(v)$ is the
reference angle. Densities are pushed through these maps. A primitive
instead composes with $U_g$: if $q_u=\partial_uQ_u$,
then $Q_v(v)=Q_u(U_g(v))$ and $q_u(u)\,du=\partial_vQ_v(v)\,dv$.
We use the same letter for these representations of a primitive.
Consequently $h_0$, the true quadratic form, and their relative
determinant do not change under a further reparametrization. The
unstretched coordinate $v$ will be used for the spectral calculation.

Use the reference space $\mathcal X$ of Section~\ref{sec:reference-law}
with these two angular variables. Compose the true $H_{nn}$ with the
maps $U_g$, subtract its diagonal reference kernel, and call the
mixed derivative of that smooth difference $\mathscr K$.
Thus $H_{\rm form}(Q,Q)=h_0(Q,Q)+(Q,\mathscr KQ)$ and
$T=I+j^*\mathscr K j$, where $j$ is inclusion in physical $L^2$.
We will prove the following spectral bounds:
\begin{equation}\label{eq:G-8}
 \begin{gathered}
 \|T-I\|_1\le CL,\quad\|T\|\le C,\quad
 A_T=(T+T^*)/2\ge cL^{-2}I,\\
 N_{A_T}(a)\le CL\sqrt a\quad(0<a<a_0).
 \end{gathered}
\end{equation}
The same gap and small-value count are required for singular values
and orthogonal compressions.

\subsection{Spectral bounds and the unequal-wall determinant}

The bounds above control the logarithm at small singular values; the
bulk symbols then determine its leading integral. We begin with the
local norm and cell estimates needed for that control.

The magnetic statements in Equation~\eqref{eq:G-8} will now be verified. The embedding \(j\) in unstretched \(L^2(dv)\) has uniformly bounded norm; \(jj^*\) is the labelwise kernel
\[
 G(v,v')=\frac1{2\pi^2}\log\frac{\sin((t_g(v)+t_g(v'))/2)}
                              {\sin(|t_g(v)-t_g(v')|/2)}.
\] The kernel $G$ has uniformly integrable logarithmic singularities and
exponential off-diagonal decay. Thus Schur's test gives a uniform bound
for $j$. We next prove the lower spectral bounds directly in these
magnetic coordinates, before computing the determinant.

For a primitive $Q=(Q_X,Q_Y)$ put $d_g(v)=A_g-|v|$ and define
\[
 \mathfrak e(Q)=\sum_g\left[
 \iint_{|v-v'|\le2}
 \frac{|Q_g(v)-Q_g(v')|^2}{|v-v'|^2}\,dv\,dv'
 +\int_{d_g(v)\le2}\frac{|Q_g(v)|^2}{d_g(v)}\,dv\right].
\]
The odd angular extension gives
\begin{equation}\label{eq:magnetic-local-norm}
 h_0(Q,Q)\le C(\mathfrak e(Q)+\|jQ\|_2^2),\qquad
 \operatorname{Re}H_{\rm form}(Q,Q)\ge c\mathfrak e(Q).
\end{equation}
For the first inequality, the angular difference kernels in physical
variables are
\[
 \frac{t'_g(v)t'_g(v')}{\sin^2((t_g(v)-t_g(v'))/2)},\qquad
 \frac{t'_g(v)t'_g(v')}{\sin^2((t_g(v)+t_g(v'))/2)}.
\]
The first is at most $C|v-v'|^{-2}$ locally and is exponentially
integrable away from the diagonal. The reflected kernel has the same
bound away from common endpoints; near an endpoint it is bounded by
the inverse square of the sum of endpoint distances. Integrating one
variable gives $C/d_g(v)$. These statements use the exponential
reference density and its relative derivative bounds, which are uniform
also after rare-center mean matching. They prove the first inequality.

For the second, push $Q_g$ to the physical $u$ coordinate and extend
by zero outside its wall. The real-energy theorem controls the Fourier
seminorm with multiplier $k^2/(1+|k|)$. This is equivalent to the local
difference seminorm with kernel $|u-u'|^{-2}$ and a fixed sufficiently
large distance cutoff, including the zero extension. Near a cap,
$U'_g(v)\asymp\min(1,d_g(v))$, and interactions with the exterior
zeros give
$U'_g(v)/(L_g-|U_g(v)|)\asymp1/d_g(v)$. When two endpoint distances
are comparable, change of variables preserves the local difference
kernel up to bounded factors. When they are not comparable, the
inequality $|Q(v)-Q(v')|^2\le2|Q(v)|^2+2|Q(v')|^2$ bounds the
missing part by the endpoint term. The mean-matching derivative is
uniformly comparable to one, so this argument also covers rare $X$.
This proves Equation~\eqref{eq:magnetic-local-norm}.

Partition each physical parameter interval into cells of lengths between
$1/2$ and one, and let $q_a$ be the mean on cell $a$. The local
seminorm bounds the sum of the cell variances and
$\sum_a|q_{a+1}-q_a|^2$. On a block of length comparable to $b\ge1$
whose total mean vanishes, discrete Poincar\'e gives
$\|Q\|_2^2\le Cb^2\mathfrak e(Q)$. The same estimate on a whole
wall, without a zero-mean condition, follows by controlling its first
and last cell means with the endpoint term. Since both interval
lengths are at most $CL$, Equation~\eqref{eq:magnetic-local-norm}
gives the global gap $cL^{-2}$.

For a threshold $a$ above a fixed multiple of $L^{-2}$, choose blocks
of length $b=c_0a^{-1/2}$ with small fixed $c_0$. Their zero-mean
conditions have codimension at most $CL/b$; on their common kernel
the Rayleigh quotient of $A_T$ is at least $a$. Min--max gives
$N_{A_T}(a)\le CL\sqrt a$. Below the global gap the count is zero;
changing constants covers the intermediate range. Compression cannot
increase this count. Finally, if a subspace satisfies
$\|TQ\|<a\|Q\|$, then
$\langle Q,A_TQ\rangle\le\|Q\|\|TQ\|<a\|Q\|^2$ on that
subspace. This gives the same count for singular values, including
compressions. We have therefore established the magnetic lower bounds
required in Equation~\eqref{eq:G-8}. The following determinant proof
supplies its norm and trace bounds, including the rare center.

\begin{lemma}[The magnetic fluctuation determinant]\label{lem:07_magnetic_analysis-determinant}
The relative Hessian satisfies \(\|T-I\|_1=O(L),\ \|T\|=O(1)\), and
\begin{equation}
 -\tfrac12\log|\det T|
 =\tfrac{7\gamma}{8}L_X+\tfrac{5\gamma}{16}(L_Y-L_X)+o(L).
 \label{eq:M-31}
\end{equation}
In the rare case, let \(T^0\) be the straight comparison with
main coordinates \(W_X(v)+i\beta,W_Y(v)\), using exactly the same
reference form \(h_0\). Then
\begin{equation}
 \|T-T^0\|=o(1),\qquad \|T-T^0\|_1=o(L).
 \label{eq:M-32}
\end{equation}
In particular all errors in \eqref{eq:M-31} are uniform along every subsequence
specified above, including \(t\uparrow1\).
\end{lemma}
\begin{proof}
\par\smallskip\noindent\textbf{Trace bounds for the straight comparison.}
In the rare case let \(T^0\) denote the straight comparison using physical main coordinates \(W_X(v)+i\beta,\ W_Y(v)\), and \emph{exactly the same} \(h_0\). In the regular case let \(T^0=T\). For \(T^0\) the mixed derivative of the true-minus-reference kernel has uniform \(C^j\) bounds of order \(C_j e^{-c|v-v'|}\), all fixed \(j\). Indeed on a regular curved main the same diagonal argument subtracts the horizontal log, the residual using the divided ratio for the curved \(\sinh\); it is uniformly regular including at caps by composition (the ratio of the two physical complex and horizontal differences before composing with \(W_X\) already extends smoothly with controlled log). Other near-diagonal entries are offset from their zero angles. At long differences the linear and sign asymptotes, including the terms from curved heights, are additive in individual coordinates in fixed order and vanish on mixed differentiation; the remaining log tails and the reference sine ratios obey the
same exponential estimates. To obtain the trace bound, partition each
interval into cells of lengths in $[1,2]$ and use their normalized
cosine bases $e_{a,n}$. Integrating twice by parts in each variable,
with the first boundary term zero and the second controlled by a
kernel derivative, gives
\[
 |\langle e_{a,n},\mathscr K e_{b,m}\rangle|
 \le \frac{Ce^{-c|a-b|}}{(1+n^2)(1+m^2)}.
\]
Summing the trace norms of these rank-one entries gives $O(L)$;
Schur's test gives an operator norm $O(1)$. Boundedness of $j$
transfers both statements to $T^0$.

\par\smallskip\noindent\textbf{The common and shoulder symbols.}
Cyclic kernel integrals then
give the word limits used below. More precisely, deep points \(|v|<L_X\), away from the center and edges by distances tending to infinity, see the common symbol \(T_\infty(k)\) there with \(p=(3-\lim\theta_\infty)\beta\) (angle 1 for the horizontal comparison in the rare case); see the main transform formulas. Indeed the reference densities in each tail are asymptotic with derivatives to a constant times \(e^{-\gamma |v|}\), and the curved height is exponentially asymptotic to \(\theta_\infty\). Deep points in \(L_X<|v|<L_Y\) see only its \(Y,Y\) entry. Their respective weights after division by \(2L\) are \(\lim L_X/L\) and the complement. Bounded-distance exceptional regions have zero limiting share, and relative-coordinate tails can be truncated uniformly by the kernel bounds. For completeness, its $2\times2$ expression, for $k>0$ and
$x=\beta k$, is
\[
 T_\infty(k)=\frac{2\cosh x-1}{\sinh(4x)\coth(3x/2)}
 \begin{pmatrix}\cosh(3x)&e^{pk}\\ e^{-pk}&\cosh(3x)\end{pmatrix};
 \qquad T_\infty(-k)=T_\infty(k)^{\mathsf t}.
\]
It is obtained by dividing the true primitive symbol $k^2\widehat H$
by the limiting reference symbol $2\pi k\coth(\pi k/(2\gamma))$.
Its determinant is independent of $p$, and, with $y=e^{-x}$,
\[
 \sqrt{\det T_\infty(k)}
 =\frac{(1-y^3)(1-y^6)}{(1+y)(1-y^8)}.
\]
The logarithmic integral is
\[
 \int_0^\infty\log\sqrt{\det T_\infty(x/\beta)}\,dx
 =-\frac{\pi^2}{18}-\frac{\pi^2}{36}
       -\frac{\pi^2}{12}+\frac{\pi^2}{48}
 =-\frac{7\pi^2}{48}.
\]
Each term follows by integrating the convergent logarithmic series.
Accounting for both signs of $k$, the determinant square, and the
common interval length $2L_X$ yields the contribution
$7\gamma L_X/8$ to $-\frac12\log|\det T|$.
The scalar shoulder symbol for \(x=\beta|k|,\ y=e^{-x}\) is
\[
 \frac{c_*(k)\cosh(3x)}{\sinh(4x)\coth(3x/2)}
 =\frac{(1-y^3)(1+y^6)}{(1+y)(1-y^8)}
\]
The integral of the logarithm of the shoulder symbol over
\(x\in(0,\infty)\) is
\[
 -\frac{\pi^2}{18}+\frac{\pi^2}{72}
 -\frac{\pi^2}{12}+\frac{\pi^2}{48}
 =-\frac{5\pi^2}{48},
\]
by expanding each logarithm in its absolutely integrable series. This gives \eqref{eq:M-31} by the same truncated-log argument once trace convergence holds also for \(T\), since \eqref{eq:G-8} was proved for \(T\).

\par\smallskip\noindent\textbf{Passing to the logarithmic determinant.}
Here is the truncated-log passage used in this calculation. For the
singular values $s_i(T)$, the small-value count gives
\[
 \sum_{s_i<\eta}\log(\eta/s_i)
 =\int_0^\eta\frac{\#\{s_i<r\}}r\,dr\le CL\sqrt\eta.
\]
Clipping $\log(T^*T)$ below $\eta^2$ therefore loses only
$O(L\sqrt\eta)$. On the bounded spectral interval, approximate the
clipped logarithm by polynomials vanishing at one, with error at most
$\epsilon|x-1|$. The trace error is at most
$\epsilon\|T^*T-I\|_1=O(\epsilon L)$. Cyclic kernel integration
proves the polynomial trace limits: truncate relative separations,
use the integrable logarithmic bound for $G$ and the exponential bounds
for $\mathscr K$, and then discard bounded neighborhoods of the center
and walls. The common interval occupies length $2L_X$ and the
shoulders length $2(L_Y-L_X)$. The same clipping estimate holds for
the displayed symbols. The common symbol has one singular value
comparable to $k^2$ near zero and its other singular value stays
positive; the scalar shoulder symbol tends to $3/8$ there. Both
symbols approach identity exponentially at infinity.
Take size to infinity, then $\epsilon\downarrow0$ and
$\eta\downarrow0$. The two logarithmic integrals give exactly the
coefficients in Equation~\eqref{eq:M-31}.

\par\smallskip\noindent\textbf{The rare-center comparison.}
It remains to prove the comparison uniformly at the rare central scale. Ordinary derivative bounds in \(v\) degenerate at that scale, so we use a dyadic decomposition. In intermediate coordinates put
\[
 \mathfrak z_X(b_X^{-1}(\xi))=\xi+i\beta+E(\xi).
\]
Then \(E,E'\) are \(o(1)\) uniformly, with fixed higher derivatives \(E^{(j)}=o(1)\,[\min(1,|\xi|+d_0)]^{1-j}\), by the graph and matching estimates. The kernel on primitives uses a mixed derivative of the difference \(K_d=H_{nn}-H_{nn}^0\) with the two sets of maps just described. Split \(X\), on \(|v|\le1\) (so \(W_X(v)=v\)), into sign-paired shells of radii \(r_j=2^{-j}\) down to \(r_J\asymp d_0\): \(r_{j+1}<|v|<r_j,\ j<J\), and the last central interval \(|v|<r_J\). In shell \(j\) encode \(Q_X\) by its oscillation from the single mean \(m_j\) over that shell, measured in \(L^2\) after rescaling the shell by \(r_j\), and keep the numbers \(m_0,\ m_j-m_{j-1}\). The means thus contribute to \(Q_X\) by these coefficients times the steps \(\mathbf1_{|v|<r_j}\). Outside this region and on \(Y\) keep \(Q\) itself in \(L^2\) on cells of lengths between 1 and 2. This joint encoding has squared norm bounded by \(C h_0(Q,Q)\). Indeed same- and adjacent-shell squared mean/oscillation differences are bounded in sum by the ordinary difference seminorm on \([-1,1]\), directly controlled there by the angular difference representation of \(h_0\); the outer mean and ordinary cells cost \(\|jQ\|^2\).
For example, with \(A_j=\{r_{j+1}<|v|<r_j\}\),
\[
 \frac1{r_j}\int_{A_j}|Q-m_j|^2
 \le C\iint_{A_j\times A_j}
        \frac{|Q(u)-Q(v)|^2}{|u-v|^2}\,du\,dv,
\]
and the same calculation on \(A_j\times A_{j+1}\) bounds
\(|m_j-m_{j+1}|^2\). These pair regions have bounded overlap.
The last central interval satisfies the same scaled estimate. Thus the
encoding map is bounded independently of the number of shells.

For oscillation or ordinary cells the matrix kernel on this encoding takes one derivative of \(K_d\) on the natural unit scale of each applicable variable (rescaled for shells). For a step, integrate that derivative over the step: take instead the difference of values at \(+r_j,-r_j\). Thus there is no cost from a large \emph{common} primitive mean at tiny scales. Each cell pair, including derivatives through two further orders per continuous variable, has bounds \(C\epsilon_L\) times:
\begin{itemize}
\item \(\min(r_i/r_j,r_j/r_i)\) between shells (steps are included at their corresponding scale);
\item \(r_j e^{-c|v'|}\) from shell \(j\) to an ordinary cell at \(v'\);
\item \(e^{-c|v-v'|}\) between ordinary cells;
\end{itemize}
where \(\epsilon_L=o(1)\).

To verify them, on two central \(X\) arguments the diagonal log change is a multiple of the smooth log of \(1+(E(v)-E(v'))/(v-v')\) plus regular analytic changes. For comparable scales this has all required bounds by divided differences at nearby arguments in the same regular neighborhood of radius of that scale (including the innermost pair), or by separation otherwise. For a much smaller scale \(r_i\ll r_j\) differentiation on the smaller scale gains \(r_i/r_j\), since the denominator has size \(r_j\) and all positive-order normalized derivatives of \(E\) on the smaller shell have size \(C\epsilon_L r_i\); further derivatives on either scale preserve the bound. A step difference gains the same by varying its argument across \([-r_i,r_i]\), using only the uniform slope bound for \(E\) there. The larger step endpoints themselves stay on the larger scale. The argument across the interface to ordinary \(X\) cells is identical; a bounded-distance ordinary \(Y\) argument instead sees a nonsingular analytic kernel. In far-separated comparisons the pure one-variable log asymptotes drop out after the two derivatives/differences, including when one uses step endpoints (order with the other far argument is fixed); the log series remainders give exponential decay along with the small change bounds, with the scale gain from a central derivative or step. On ordinary \(X\) cells near a cap the divided differences for the change use \(E(\xi)-E(\xi')\) divided by \(\xi-\xi'\) \emph{before} the cap map, so do not incur any endpoint derivative loss. These observations prove the stated estimates for the log branches continued by the graph deformation; additive constants do not enter.

By smooth cell expansion (e.g. the same cosine trace bound as before, on the two separate sign components where needed), these are also trace-norm bounds per pair. Their row and column sums are \(O(\epsilon_L)\), and their overall sum \(O(\epsilon_L(L+J))\). Factoring on both sides by the bounded encoding proves \eqref{eq:M-32}. Consequently the needed trace bounds and word limits transfer to \(T\) (trace-norm comparison of products), proving \eqref{eq:M-31}.


Since \(J=O(1+\log(1/d_0))=O(L)\), the shell-pair trace estimate is
indeed \(o(L)\). To transfer a word of fixed length in \(T-I,T^*-I\),
expand its difference from the corresponding word in \(T^0\) as a
sum with one difference factor in each summand. The trace norm is
\(o(L)\), because all other factors have bounded operator norms.
The small-singular-value count in Equation~\eqref{eq:G-8} controls the
removed portion of the logarithm uniformly, so no uniform pointwise
inverse bound at the rare scale is required.

\end{proof}

\subsection{A finite core and regular coordinates}

The determinant estimate controls the global cost. To realize the trial
field by moving actual particles, we next choose a finite core and a
coordinate in which its fields, dual tests, and density have uniform
regularity.

Fix a small tolerance $\xi>0$. We construct a real orthogonal
projection $\Pi$ of rank $O_\xi(L)$, supported on each label in
$|v|<(1-\delta_*)L_g$ for a fixed $\delta_*>0$, such that
\begin{equation}\label{eq:G-10}
 \|T-I-\Pi(T-I)\Pi\|_1\le\xi L.
\end{equation}
Apply the smooth-core construction of \cite[Lemma~3.11]{Electric}
to $T^0$. Its cell approximation uses the compensated-kernel bounds
proved above. It then applies the reference inverse $G$ to smooth
cell functions and cuts off the resulting fields inside the walls.
The exponential decay of $G$ makes this cutoff error negligible;
discarding the outer $O(\delta_*L)$ cells costs $O(\delta_*L)$ in
trace norm. These arguments apply separately on our two intervals.
On rare $X$, both $G$ and the reference density in $v$ remain regular,
so the constructed orthonormal fields have polynomial fixed-order
derivative bounds in $v$. To pass from $T^0$ to $T$, use
\[
 \|(T-I)-\Pi(T-I)\Pi\|_1
 \le \|(T^0-I)-\Pi(T^0-I)\Pi\|_1+2\|T-T^0\|_1.
\]
Choosing the initial tolerance smaller and using
Equation~\eqref{eq:M-32} proves Equation~\eqref{eq:G-10}.

Write $S=\operatorname{ran}\Pi$, $R=S^\perp$ and
$T_e=T_{RR}-T_{RS}T_{SS}^{-1}T_{SR}$. Its Hermitian part
$\operatorname{Re}T_e$ inherits the gap and small-value count in
Equation~\eqref{eq:G-8}. The accompanying comparison is
\begin{equation}\label{eq:G-11}
 \begin{aligned}
 |\log|\det T_{SS}|-\log|\det T||&=o_\xi(L),\\
 \sum_{\lambda<1,\ \lambda\in\operatorname{spec}\operatorname{Re}T_e}
             -\log\lambda&=o_\xi(L).
 \end{aligned}
\end{equation}
Here the normalized limsup tends to zero when $\xi\downarrow0$
after $L\to\infty$. These conclusions follow from
\cite[Lemma~3.12]{Electric}, whose operator argument uses precisely
Equation~\eqref{eq:G-10} and the magnetic gap and small-value counts
already proved. In particular it controls the negative logarithms of
$\operatorname{Re}T_e$, not only the determinant of $T_e$.

The core fields themselves are now regular in $v$. Their images under
the true form can still vary on the rare scale. The next change of
coordinate controls these dual tests as well as the particle motion;
it preserves the forms and the projection just constructed.


\begin{lemma}[A regular parameter for transport]\label{lem:07_magnetic_analysis-stretch}
On rare $X$, define $\tau(0)=0$ and
\[
 \frac{d\tau}{dv}=1+\frac1{\sqrt{d_0^2+W_X(v)^2}};
\]
on the other intervals set $\tau=v$. Let $\widetilde A_g$ be the
new endpoints and $\lambda_g(\tau)=\lambda_g^0(v)\,dv/d\tau$ the
reference probability density. Then $\widetilde A_g\asymp L$,
$|\log\lambda_g|=O(L)$, and $\lambda_g$ has smooth even endpoint
reflections with uniformly bounded relative derivatives of every
fixed order. Off the bounded-width caps, the maps extend analytically
to fixed small disks in $\tau$. In those disks the density has relative
bounds, and increments and positive-order derivatives of $v$, $U_g$,
and $\mathfrak z_g(U_g)$ are bounded on the scale $dv/d\tau$.
For every fixed $\delta_*>0$, if $N_g\asymp s_g$, then
\begin{equation}\label{eq:M-33}
 N_g\lambda_g(\tau)\ge e^{c(\delta_*)L}
 \qquad\text{on }|v|\le(1-\delta_*)L_g
\end{equation}
for all sufficiently large sizes.

Suppose also that $n_e,\ell_e,|M|$ are bounded by fixed powers of $L$.
In the $\tau$ coordinate the following data have polynomial bounds
for every fixed derivative order:
\begin{itemize}
\item the compensated kernel $H_{nn}-H_{\rm ref}$, with smooth even
endpoint reflections and analytic continuations in its core variables;
\item the linear density test
$H_{nn}(m'p_L-\mu)+\widetilde H_{ne}\zeta-\mathfrak t_n$, with the
same reflections and core analyticity;
\item the smooth dual tests obtained by applying $h_0$ or
$H_{\rm form}$ to the core basis fields or the trial primitives.
These tests have odd reflections when paired against primitives.
\end{itemize}
The constant term of $\operatorname{Re}\mathcal D$ on $m'p_L$ is
bounded below by $-{\rm poly}(L)$. The core density bound also holds
where a trial primitive is nonzero through the center.
\end{lemma}
\begin{proof}
\par\smallskip\noindent\textbf{Local scale and density.}
On rare $X$,
\[
 \frac{dv}{d\tau}\asymp\min(1,d_0+|v|).
\]
The added length near the center is
$O(1+\log(1/d_0))=O(L)$. Thus the intervals still have length
comparable to $L$. In a $v$-disk of radius $c(|v|+d_0)$ near the
center, the derivative of $\tau$ has small relative variation when
$c$ is small. Analytic inversion, together with the graph and
mean-matching estimates on that same scale, gives the claimed fixed
$\tau$-disks. At the caps the stretch is smooth and respects the
reflections; analytic continuation there is not needed.

The identity $\lambda_g=\lambda_g^0\,dv/d\tau$ gives the logarithmic
and relative derivative bounds. On $|v|\le1$ in rare $X$ it gives
\[
 N_X\lambda_X\ge c s d_0\ge e^{c'L},
\]
using the minority margin. Away from this central region, throughout
a core,
\[
 N_g\lambda_g\ge c s_g e^{-\gamma(1-\delta_*)L_g}
       =c e^{\gamma\delta_*L_g}.
\]
These two estimates prove Equation~\eqref{eq:M-33}. In particular
one need not remove the center from the support of a primitive.

\par\smallskip\noindent\textbf{Compensated logarithms.}
Consider first two points on the same label. If their $\tau$-distance
is at least a fixed $c_0>0$, their horizontal physical separation is
at least a constant times either local scale $dv/d\tau$. This follows
by integrating the local scale on a fixed subinterval; it remains true
at a cap, where the physical derivative can vanish at the endpoint.
Likewise the sine of half the angular difference bounds both local
$\lambda_g$ from below up to constants. The sine of the angular
half-sum is at least as large. The logarithms and their derivatives
therefore have the asserted bounds away from the local diagonal.

Near the diagonal, divide both physical and reference differences by
$\tau-\tau'$. Off the caps, the physical complex quotient has size
comparable to $dv/d\tau$, while the reference cosine quotient has
size comparable to $\lambda_g\sin t_g$; here
$\sin t_g\asymp\lambda_g^0$ away from the endpoints. Their relative
analytic derivative bounds follow from divided differences on the
fixed disks just constructed. At an endpoint, divide also by the
reflected coordinate difference. The quadratic cap maps make both
quotients smooth, nonzero on their natural scales, and even under
reflection. Their logarithmic ratio is consequently smooth through
both diagonals. Cross-label factors and the remaining same-label
factors stay at regular imaginary offsets when their real arguments
approach. Shrinking the core disks preserves those separations.
This proves the bounds and core analyticity for
$H_{nn}-H_{\rm ref}$.

\par\smallskip\noindent\textbf{Linear sources and end corrections.}
For each label decompose the signed measure $m'_gp_g-\mu_g$ into
$(m'_gC_g-s_g)\bar\nu_g\,du$ plus $m'_gC_g$ times the difference
between the unnormalized capped push and $\bar\nu_g\,du$.
The first coefficient is polynomial. Its convolution is given by
the mean identities, so it consists of ordinary $V$-logs evaluated
on the contour. The only possible shrinking site-pole distance is
at the rare $X$ center. There it is bounded below by a constant times
$|u|+d_0$, exactly the local scale controlled by $\tau$.

For the difference of measures, insert a fixed-width cutoff around
the caps and their exterior tails. The measures agree off this set,
and their densities in their respective parameters carry the factor
$e^{-\gamma L_g}$. In a same-label end convolution, subtract the
ordinary local logarithm and its reflected term for the capped input.
The curved divided-log remainder is regular. Derivatives of the
singular logarithm transfer to the smooth cutoff density, using its
even extension for the capped input; for the full-line input transfer
derivatives in $u$ before composition. Multiplication by $m'_g$ thus
leaves polynomial bounds. The end correction is analytic in core
variables by spatial separation on the same label and by imaginary
offsets on the other label.

The explicit-particle term $\widetilde H_{ne}\zeta$ uses the regular
offsets and horizontal subtraction from Equation~\eqref{eq:M-21}.
Its subtraction $R_h$ is inactive on moving core variables. The anomaly
logs have the uniform offsets established in the mean calculation.
These observations prove the linear-test assertion. Its constant
term has the stated lower bound by Equation~\eqref{eq:M-23}.

\par\smallskip\noindent\textbf{Dual tests and polynomial approximation.}
Against a smooth density supported in the core, subtract the local
$\tau$-diagonal logarithm of $H_{\rm ref}$ and transfer its derivatives
to the density. The remaining kernel has the bounds just proved.
Thus $H_{\rm ref}$ maps these densities to smooth even tests with
polynomial derivatives. The basis fields were constructed with
polynomial derivatives in $v$, hence also in $\tau$. The trial
primitives have the same property by Equations~\eqref{eq:M-24},
\eqref{eq:M-26} and~\eqref{eq:M-27}. Differentiating their density
tests gives the odd-reflected primitive tests for $h_0$. Adding the
smooth compensated kernel gives those for the true $H_{\rm form}$.
This last step is why no polynomial derivative bound for the true
form in the unstretched rare coordinate is needed.

We will also use physical sine approximation on the $\tau$ intervals.
Two estimates ensure that its errors remain polynomial. First,
$\mathcal X\hookrightarrow L^2(d\tau)$ has polynomial norm:
the inverse kernel satisfies
\[
 |G(\tau,\tau')|\le C\bigl(L+
       |\log\min(1,|\tau-\tau'|)|\bigr)
\]
by the exponential lower bound on $\lambda_g$, and Schur's test
applies on intervals of length $O(L)$. Second, $h_0$ is polynomially
bounded on unit vectors in every polynomial Dirichlet sine bandwidth.
For this use the angular difference kernel with its density factors.
Its unreflected part is at most
$C(1+|\tau-\tau'|^{-2})$, by the relative density derivative bounds.
The reflected part contributes an endpoint inverse-distance mass term,
since
\[
 \sin t_g\ge c\lambda_g\min(1,\widetilde A_g-|\tau|),
\]
and its density-weighted kernel is bounded outside pairs near the same
endpoint. Sine functions vanish
at the endpoints, so these local difference and boundary estimates
give the asserted bandwidth bound.
\end{proof}

\subsection{Source and compensated quadratic estimates}\label{sec:reference-integration}

The preceding construction gives the concrete magnetic coordinates and
trial form. We now state the two reference-law bounds used to integrate
that form, retaining the partition masses from
Equation~\eqref{eq:reference-normalization}.

\begin{proposition}[Reference sources]\label{prop:reference-source}
For fixed $C_1$, integer dimensions $0\le N_g\le C_1e^{\gamma L}$,
and a real smooth reflected test $\ell(Q)$,
\begin{equation}\label{eq:G-13}
 \log\mathbb E_0 e^{\ell(Q)}\le C_2L+\tfrac12\|\ell\|_*^2.
\end{equation}
If $\ell(Q)=(f,q)$ and
$f_g(t)=f_{g,0}+2\sum_{k\ge1}f_{g,k}\cos kt$, then
$\|\ell\|_*^2=\sum_{g,k\ge1}k|f_{g,k}|^2$.
Constants in $f$ do not affect the test.
\end{proposition}

This is the reference-source estimate of
\cite[Theorem~4.2]{Electric}: comparison to the reflected circle
integral gives the displayed Gaussian term, and the reciprocal factors
$J_{N_g}^{-1}$ give the $O(L)$ normalization cost. The circle inequality
is a finite strong Szeg\H{o} bound \cite{Simon2005Szego}; its
Schur-polynomial proof is related to the unitary trace-moment argument
of Diaconis and Evans \cite{DiaconisEvans2001}.

For the sharp bound we need to remove this $O(L)$ normalization cost.
Set
\[
 \mathcal G=\{\|Q\|_{-2,\widetilde A}\le L^D\},
\]
where the negative norm uses squared weights
$(1+(n/\widetilde A_g)^2)^{-2}$ in the orthonormal sine basis of
each $\tau$ interval. The event may couple the two labels.
For fixed $D$, dimensions $N_g\asymp s_g$, and smooth even-reflected
physical tests $f_g(\tau)$ whose fixed-order derivatives have specified
polynomial bounds, we claim
\begin{equation}\label{eq:G-14}
 \log\mathbb E_0[\mathbf1_{\mathcal G}e^{\ell(Q)}]
 \le o(L)+\tfrac12\|\ell\|_*^2.
\end{equation}
The error is uniform when these powers and dimension-comparison
constants are fixed. Theorem~4.3 of \cite{Electric} proves the
corresponding electric statement. We adapt its change of variables
here, since the unequal walls and the rare center affect both the
motion and the normalization gain.

Choose a smooth primitive $0\le a_g\le1/4$, equal to $1/4$ on
$|v|\le(1-2\delta)L_g$ and zero before
$|v|=(1-\delta)L_g$, with monotone slow switches in $v$. For the two signs
$\epsilon=\pm1$ make the real changes of variables
\[
 F_{g,\epsilon}(\tau)
       =\tau-\epsilon\frac{a_g(\tau)}{N_g\lambda_g(\tau)}.
\]
These are increasing diffeomorphisms, equal to the identity near the
endpoints. Their displacements and all required fixed derivatives are
exponentially small by Equation~\eqref{eq:M-33}, including across the
center. The cosine-difference ratios extend smoothly through the
same-label diagonal by analytic division. Off that diagonal the sine
separation estimates in Lemma~\ref{lem:07_magnetic_analysis-stretch}
give the same exponentially small derivative bounds. Thus their errors
paired with $Q$ or $Q\otimes Q$ are $o(1)$ on $\mathcal G$.
Convolutions of $H_{\rm ref}$ with the smooth mean-push remainders
have the same bounds by local-log subtraction. These are the error
estimates used in the proof of \cite[Theorem~4.3]{Electric}, now in
$\tau$.

Write $\mathfrak g_g=\log\cot(t_g/2)$ and
\[
 G_g=(\mathfrak g_g,a_g')-\tfrac12h_0(a_g,a_g),
\]
with the form restricted to label $g$. The one-coordinate density
ratio cancels the omitted half-diagonal in the Vandermonde log change.
As in that proof, pulling back the tilted density with angular weight
$2(1+\epsilon\cos t_g)$ instead of $2\sin t_g$ gives its original
density times
\[
 \exp\left\{\epsilon\bigl[(\mathfrak g_g,X_{N_g})
          +(f_g,a_g')-(q_g,H_{\rm ref}a_g')\bigr]+G_g+o(1)\right\}
\]
on $\mathcal G$. Only core derivatives of $\mathfrak g_g$ occur,
since the motion vanishes near the endpoints.

The gain is most simply calculated in the unstretched $v$ coordinate.
There $-\mathfrak g_g'=t_g'/\sin t_g\to\gamma$ in either tail,
so
\[
 (\mathfrak g_g,a_g')=\gamma\int a_g(v)\,dv+o_\delta(L).
\]
On the support of the switches, two points in the same tail satisfy
\[
 H_{\rm ref}(t_g(v),t_g(v'))
 =4\gamma\min(|v|,|v'|)
   +O\bigl(1+|\log\min(1,|v-v'|)|\bigr),
\]
whereas opposite-tail values are bounded. This follows by factoring
the difference of the squared exponentially small distances to the
nearest angular endpoint. The identity
\[
 \iint_{[0,\infty)^2}\min(r,r')b'(r)b'(r')\,dr\,dr'
       =\int_0^\infty b(r)^2\,dr
\]
and the slow switches therefore give
\[
 h_0(a_g,a_g)=4\gamma\int a_g(v)^2\,dv+O_\delta(1),
 \qquad
 G_g\ge\frac{\gamma L_g}{4}-O(\delta L)-o_\delta(L).
\]
In particular the extra length introduced by the central stretch
has no effect on this coordinate-invariant energy.

Use the \emph{same} sign $\epsilon$ for both labels before taking
the geometric mean. The terms odd in $\epsilon$ cancel on the single
event $\mathcal G$. Cauchy--Schwarz, removal of that event from the
two positive integrals, and the circle comparison then give
\[
 \log\mathbb E_0[\mathbf1_{\mathcal G}e^{\ell(Q)}]
 \le-\sum_gG_g-\sum_g\log J_{N_g}
       +\tfrac12\|\ell\|_*^2+o(1).
\]
Since $-\log J_{N_g}=\gamma L_g/4+O(1)$, the excess over the
Gaussian term is $O(\delta L)+o_\delta(L)$. Taking $L\to\infty$
and then $\delta\downarrow0$ proves Equation~\eqref{eq:G-14}.

\begin{proposition}[Compensated quadratic integration]
\label{prop:quadratic-integration}
Let $F$ be real symmetric and bounded on $\mathcal X$, with
$F\ge gI$, $g>0$, and $F-I$ trace class. Suppose its perturbation
and a real linear functional $h_0(l,\cdot)$ are represented by smooth
reflected physical kernels and tests. Let $c\in\mathbb R$. For
\[
 V(Q)=c+h_0(l,Q)+\tfrac12h_0(Q,(F-I)Q),\qquad
 m_F=c-\tfrac12h_0(l,F^{-1}l),
\]
one has
\begin{equation}\label{eq:G-18}
 \log\mathbb E_0e^{-V(Q)}\le -m_F+C_2L
        +\tfrac12\sum_{\lambda<1,\ \lambda\in\operatorname{spec}F}
                   -\log\lambda.
\end{equation}
The sum counts multiplicity. If the sharper source bound
\eqref{eq:G-14} holds, the gap is inverse-polynomial, and the data
and required fixed-order derivatives have polynomial bounds, the
same estimate on $\mathcal G$ has $o(L)$ in place of $C_2L$.
\end{proposition}

For clarity, the cancellation underlying this import is short.
Set $Q_*=-F^{-1}l$ and $J=(I-F)_+$. Completing the square gives
\[
 -V(Q)\le-m_F+h_0(Q_*,Q)-\tfrac12\|Q_*\|^2
              +\tfrac12h_0(Q-Q_*,J(Q-Q_*)).
\]
Linearize the last term with a real Gaussian vector $g_J$ of covariance
$J$, and apply Proposition~\ref{prop:reference-source} to
$Q_*+g_J$. The mixed terms cancel exactly, leaving
$e^{-m_F+C_2L}\mathbb E e^{\|g_J\|^2/2}
=e^{-m_F+C_2L}\det(I-J)^{-1/2}$. Smooth finite-dimensional
approximation, gap control and Fatou's lemma give the stated general
form, as in \cite[Lemma~4.4]{Electric}. For its sharpened version,
use physical sine truncation in $\tau$. The two Sobolev estimates at
the end of Lemma~\ref{lem:07_magnetic_analysis-stretch} give polynomial
operator and dual-test bounds on every polynomial bandwidth. Smoothness
then makes the truncation error smaller than any prescribed inverse
power of $L$, both on $\mathcal G$ and after division by the gap.
Equation~\eqref{eq:G-14} applies to the resulting polynomially bounded
Gaussian sources; outside a sufficiently large polynomial ball use
Equation~\eqref{eq:G-13} and the Gaussian tail. This is the approximation
argument of that lemma with our physical coordinate.

\subsection{The finite-rank transport construction}

The remaining step is to realize the complex trial field in the
particle integral. We adapt the three-stage deformation proved in
\cite[Lemma~5.2]{Electric}. Its statement concerns electric sectors;
below we give the affine path for the magnetic form, prove its gap
bounds, and check the geometric estimates needed for the same contour
argument. Particle dimensions stay fixed throughout.

For fixed explicit configuration $\zeta$, define the quadratic cost
$\mathscr C_0(Y)=\mathcal D(m'p_L+Y',\zeta)$, where density
 derivatives are pushed from the regular coordinate to the physical one.
For a real smooth trial primitive $Z_0$ supported in the active core,
write
\[
 C_z(Y)=\mathscr C_0(Y+iZ_0)
       =c_z+h_0(u_z,Y)+\tfrac12h_0(Y,TY),\qquad
 m_z=\inf_{Y\ {\rm real}}\operatorname{Re}C_z(Y).
\]
The endpoint of the affine transformation is
\begin{equation}\label{eq:G-19}
 Q+a(Q)=iZ_0+Q_R+T_{SS}^{-1/2}Q_S
       -T_{SS}^{-1}\Pi u_z-T_{SS}^{-1}T_{SR}Q_R.
\end{equation}
The square root is the principal branch in the right half-plane,
well-defined by accretivity. The corresponding motion of a coordinate
$r$ of label $g$ is $-a_g(Q)(r)/(N_g\lambda_g(r))$.
Throughout this motion each $N_g$ is a fixed nonnegative integer.

The following proposition bounds the resulting magnetic integral,
including its two unnormalized reference masses.

\begin{proposition}[Implementation of magnetic trial actions]\label{prop:07_magnetic_analysis-transport}
Let \(Z_0\) be the primitive of the coarse or sharp tilt, and set
\[
 C_z(Y)=\mathcal D(m'p_L+(Y+iZ_0)',\zeta),\qquad
 m_z=\inf_{Y\in\mathcal X\text{ real}}\operatorname{Re}C_z(Y).
\]
Derivatives here denote densities pushed from \(\tau\) to the physical
coordinate. After the reference transport, the logarithm of the modulus
of the unnormalized sector integral is at most
\[
 -m_z+\tfrac{7\gamma}{8}L_X+\tfrac{5\gamma}{16}(L_Y-L_X)
       -\tfrac{\gamma}{4}(L_X+L_Y)+o_{\rm proj}(L)+o(L).
\]
The discarded part can be made smaller than
\(\exp(-L^{E_2})\) for any required fixed \(E_2\). The coefficient of
\(o_{\rm proj}(L)\) is arbitrarily small after the core projection is
improved. Without a tilt one also has the bound \(-m_0+O(L)\),
which permits polynomial truncation of \(n_e,|M|,\ell_e\).
\end{proposition}
\begin{proof}
\par\smallskip\noindent\textbf{Coarse integration and count truncation.}
For any chosen tilt, the infimum $m_z$ satisfies its deterministic
trial lower bound: first test smooth real fields, then use continuity
in $\mathcal X$. For $Z_0=0$, take absolute values and apply
Proposition~\ref{prop:quadratic-integration} with
$F=\operatorname{Re}T$. The smooth compensated data are defined at
fixed sizes even before polynomial truncation. The spectral bounds
and $\|T-I\|_1=O(L)$ give an $O(L)$ negative-log sum, and the
reference masses $J_{N_g}$ are bounded above. Thus the logarithm of
the unshifted integral is at most $-m_0+O(L)$.

Combining this with Equations~\eqref{eq:M-20} and~\eqref{eq:M-23}
retains damping in
$\ell_e+(M^2+n_e^2)/L$, apart from $C_{h,\varepsilon}(L+n_e)$.
Sum the Gaussian damping in $M$ and the position damping with the
factorial measures as in Equation~\eqref{eq:M-30}. We may thereby
restrict $n_e,|M|,\ell_e$ to sufficiently large fixed powers of $L$
with error $\exp(-\omega(L))$ at the extracted-product normalization.
Fix these powers from now on. Then $N_g=s_g+O(L^K)\asymp s_g$,
and all kernels, dual tests, and lower bounds for constant terms have
the polynomial controls proved above. The sharp source estimate
\eqref{eq:G-14} is available in $\tau$.

\par\smallskip\noindent\textbf{The affine path and its real gap.}
Put
\[
 \mathsf A=T_{SS},\qquad \mathsf C=T_{SR},\qquad
 \mathsf B=\mathsf A^{-1}\mathsf C,
 \qquad H_\vartheta=(1-\vartheta)I+\vartheta\mathsf A,
 \quad D_\vartheta=H_\vartheta^{-1/2}.
\]
For $0\le\vartheta\le1$, the three successive output primitives are
\[
\begin{aligned}
 &Q+i\vartheta Z_0,\\
 &iZ_0+Q_R+Q_S
       -\vartheta(\mathsf A^{-1}\Pi u_z+\mathsf BQ_R),\\
 &iZ_0+Q_R+D_\vartheta Q_S
       -\mathsf A^{-1}\Pi u_z-\mathsf BQ_R.
\end{aligned}
\]
The stages join at their endpoints, and the last endpoint is
Equation~\eqref{eq:G-19}. Accretivity places the spectrum of
$H_\vartheta$ in the right half-plane. The resolvent representation
of its principal inverse square root gives polynomial bounds on
$D_\vartheta$ and its required parameter derivatives, using
Equation~\eqref{eq:G-8}. Flattening the parameter at the joins
produces a smooth path.

We verify the real Hessian along this path. Write
$T[Y,Y]=h_0(Y,TY)$ with the complex bilinear convention.
The first stage changes only the source. For real $s\in S,r\in R$,
the second-stage linear map satisfies the exact identity
\[
 T[(s-\vartheta\mathsf Br,r),(s-\vartheta\mathsf Br,r)]
 =T[(s,(1-\vartheta)r),(s,(1-\vartheta)r)]
       +(2\vartheta-\vartheta^2)T_e[r,r].
\]
Both real forms on the right have the magnetic gap, and the two
coefficients of $\|r\|^2$ add to one. At the third stage the
Hessian is block diagonal, with blocks $T_e$ and
$\mathsf A H_\vartheta^{-1}$, because $D_\vartheta$ commutes with
$\mathsf A$ and is symmetric. The Hermitian identity
\[
 H_\vartheta^*\operatorname{Re}(\mathsf A H_\vartheta^{-1})
       H_\vartheta
 =(1-\vartheta)\operatorname{Re}\mathsf A
       +\vartheta\mathsf A^*\mathsf A
\]
gives an inverse-polynomial gap throughout this stage as well.
All nontrivial input uses of $Q$ are through smooth dual tests; all
output changes lie in $S+\mathbb C Z_0$ and are core-supported.
Their polynomial bounds follow from the core construction and the
stretched regularity lemma. The path changes the Hessian by finite
rank $O(L)$, so its trace norms and lower bounds for minimizing costs
are polynomial too.

At the endpoint, the real Hessian is
$I_S\oplus\operatorname{Re}T_e$. Its real infimum is at least $m_z$.
Indeed, for fixed real $r$, the complex stationary point of
$C_z(s+r)$ in $S$ is
$s_*=-\mathsf A^{-1}(\Pi u_z+\mathsf Cr)=s_1+is_2$. Completing
that square gives
\[
 \operatorname{Re}C_z(s_*+r)
 =\operatorname{Re}C_z(s_1+r)
       +\tfrac12h_0(s_2,\operatorname{Re}\mathsf A\,s_2)
 \ge m_z.
\]
The remaining endpoint input adds $\|s\|^2/2$.

\par\smallskip\noindent\textbf{The particle chain and its pole distances.}
Write $a_\vartheta(Q)$ for the output change along the flattened
path. Choose a smooth cutoff $\chi$, equal to one on $[0,1]$ and
zero on $[4,\infty)$, and use the parameter
\[
 \vartheta(Q)=\chi(\|Q\|_{-2,\widetilde A}^2/L^{2D}).
\]
Thus configurations in $\mathcal G$ undergo the full path, while
those outside $\|Q\|_{-2,\widetilde A}\le2L^D$ remain fixed.
For $a(Q)=a_{\vartheta(Q)}(Q)$, move a particle at $\tau$ by
\[
 \Delta_g(\tau;Q)=-\frac{a_g(Q)(\tau)}{N_g\lambda_g(\tau)}.
\]
The homotopy uses $a_{\upsilon\vartheta(Q)}(Q)$ for
$0\le\upsilon\le1$.
The exponent $D$ will be chosen after the polynomial truncations.
For fixed $D$, Equation~\eqref{eq:M-33} makes $\Delta_g$ and its
fixed-order spatial derivatives exponentially small, with $Q$ and
the cutoff parameter frozen for these derivatives.

On rare $X$, a fixed small $\tau$-disk corresponds to physical
increments on scale $dv/d\tau\asymp\min(1,d_0+|u|)$ in the core.
Hence the physical particle displacement is exponentially small
relative to $d_0+|u|$ near the center. The nearest site pole has at
least that distance by the stretched regularity lemma. All other
poles have fixed separation there. The compensated kernels and maps
therefore remain analytic on the moving chains, including against
inactive coordinates. Same-label coincidences are Vandermonde zeros;
the divided difference of the compensated ratio extends analytically
through them.

The exact contour argument of \cite[Lemma~5.2]{Electric} now applies
to these chains. Its maps need only be $C^2$ in the real particle
variables. This holds here because a particle's first and second
variations of its empirical primitive are $-\delta_\tau$ and
$\delta_\tau'$, both continuous in the physical $H^{-2}$ space.
Holomorphy is required of the density times volume form, not of the
smooth core fields defining the chain.

To see that the one nonholomorphic factor causes no difficulty,
partition every coordinate into a slightly enlarged active core and
its complement, and freeze the complementary coordinates. A like-main
explicit particle lies outside its own wall, whereas $R_h$ is
supported at horizontal distance at most one. Thus
$A=\exp(-(X_{\rm int},R_h\zeta))$ is independent of every active
coordinate. Factor it out before applying Stokes' theorem to the
remaining holomorphic top form. The displacements vanish near the
core boundaries, so the side faces of the homotopy contribute zero.
Only after this contour equality do we use $0\le A\le1$.

\par\smallskip\noindent\textbf{Compensated cost and finite Jacobian.}
We indicate why the error estimates in that contour proof persist
in $\tau$. Freeze $Q$ and $a$, and multiply $\Delta$ by a scalar
complex parameter. Its analytic disk can have exponentially large
radius: the motion stays inside the fixed $\tau$-disks, and the
cosine-difference ratios stay near one by simultaneous divided
differences and the small spatial derivatives of $\Delta$.
After removal of the unchanged explicit-pair constant, all measures
have at most exponential total variation and all remaining kernels
have polynomial bounds. Cauchy's estimate on the growing disk,
at a sufficiently high fixed Taylor order, therefore gives an
$o(1)$ remainder at scalar parameter one.

For real frozen displacements the $k$th mean-push coefficient is
\[
 d_{g,k}=\frac{(-\partial_\tau)^k}{k!}
              (N_g\lambda_g\Delta_g^k).
\]
It has $d_{g,1}=a_g'$, while the coefficients for $k\ge2$ and their
fixed derivatives are exponentially small, since
$N_g\lambda_g\Delta_g^k=-a_g\Delta_g^{k-1}$.
Pairings of the centered push terms against smooth tests are also
exponentially small on the polynomial event. Local-log subtraction
controls $H_{\rm ref}d_{g,k}$. The resulting finite Taylor identities
are polynomial in the frozen field coefficients and hence hold for
complex displacements as well. This is exactly the compensated
substitution in the companion proof: after its one-coordinate density
and diagonal cancellation, the cost is
\[
 \mathscr C_0(Q+a(Q))-\tfrac12h_0(Q,Q)+o(1).
\]
As before, the displayed expression means its smooth compensated
expansion; no empirical self energy is evaluated separately.

The remaining Jacobian differs from the identity by a matrix of rank
$O(L)$. On the full-shift event its reduced entries are empirical
quadratures of smooth core tests divided by $N_g\lambda_g$.
Their centered errors are exponentially small by
Equation~\eqref{eq:M-33} and the $H^{-2}$ event bound. The matrix
determinant lemma therefore gives its modulus as
\[
 (1+o(1))|\det T_{SS}|^{-1/2},
\]
just as in \cite[Lemma~5.2]{Electric}.

\par\smallskip\noindent\textbf{The transition and restoration of masses.}
For each fixed path parameter the compensated real cost has the
inverse-polynomial gap and polynomial bounds proved above. The
negative-metric penalty $L^{-E}\|Q\|_{-2,\widetilde A}^2$ has
polynomial operator norm and trace norm by the stretched Sobolev
estimates. Choose $E$ large enough that subtracting it preserves
half the gap. For this metric penalty use the sine-truncation extension
of the coarse quadratic bound in \cite[Lemma~4.4]{Electric}; its finite
smooth truncations converge on empirical fields, and Fatou's lemma
passes the bound to the full penalty. It then suppresses
$\|Q\|_{-2,\widetilde A}>L^D$ by
\[
 \exp(-L^{2D-E}+L^{E_1}),
\]
where $E_1$ is fixed independently of $D$.

The actual parameter depends on $Q$. On the transition region,
path derivatives of the cost are at most $L^{C+2D}$, so a grid of
mesh $L^{-C-2D-1}$ compares every cost with a grid cost within one.
Summing the preceding estimate over this polynomial grid controls
that dependence. The cutoff adds one rank-one Jacobian term; the
reduced matrix has size $O(L)$ and entries at most $e^{C_DL}$,
so its determinant costs at most $\exp(O_D(L^2))$. These estimates
give a discarded contribution smaller than any prescribed
$\exp(-L^{E_2})$ by choosing $D$ sufficiently large. Choose $E_2$
after the fixed polynomial count and position cutoffs: their total
factorial integration volume and number of sectors are at most
$\exp({\rm poly}(L))$, so summing the discarded contributions
remains negligible. All these choices are fixed before $L\to\infty$.

On the full-shift event apply the sharp quadratic bound using
Equation~\eqref{eq:G-14} and the infimum comparison with $m_z$.
Equation~\eqref{eq:G-11} bounds the remaining negative-log sum by
$o_\xi(L)$ and replaces the core determinant by the full determinant.
Finally restore $J_{N_X}J_{N_Y}$, whose logarithm is
$-\gamma(L_X+L_Y)/4+O(1)$. Equation~\eqref{eq:M-31} now gives
exactly the stated sector bound. Each $N_g$ remained the same integer
throughout the deformation.
\end{proof}

\subsection{The sharp bound and the polygon moment}

All factors are now available: the trial gain, the fluctuation determinant,
the two reference masses, the explicit-particle sum, and the uniform
vacuum denominator. We assemble them before summing over planar marks.

\begin{theorem}[The sharp marked-pair bound]\label{thm:07_magnetic_analysis-moment}
Fix $c_1\in(0,1]$ and a sufficiently small $c_4>0$. Consider the
physical cyclic row $[p,X,p',Y]$ with $s=|X|$, $b=|Y|$, and
$N=s+b+2$ even. All real site offsets vanish. Ordinary $Y$ sites
have height zero; ordinary $X$ sites and both marks have heights in
$\{0,\beta\}$. Let $t$ be the upper-height fraction in $X$ and set
$L_X=\gamma^{-1}\log s$, $L_Y=\gamma^{-1}\log b$.
Uniformly as $L_Y\to\infty$ under
\[
 c_1L_Y\le L_X\le L_Y,\qquad \tfrac12\le t\le1,
 \qquad t<1\ \Longrightarrow\ s(1-t)\ge e^{c_4L_Y},
\]
one has
\[
 \log\mathcal C(p,p')
 \le-\frac43\log s+\frac{617}{1200}\log(b/s)+o(\log b).
\]
\end{theorem}
\begin{proof}
For \(U>U_0\), Equation~\eqref{eq:M-24}, the count bound, and factorial
summation make the contribution negligible once \(U_0\) is sufficiently
large, after the smoothing and probe parameters have been fixed.
For \(U\le U_0\), combine Proposition~\ref{prop:07_magnetic_analysis-transport}
with Equations~\eqref{eq:M-29} and~\eqref{eq:M-30}, and divide by the
cleared denominator of Theorem~\ref{thm:06-denominator}. The denominator estimate
applies uniformly to the site mixture, and \(N\asymp b\).

The exponent can be checked before changing variables. The trial cost
contributes \(-2L_X\), the fluctuation and reference factors contribute
\(3\gamma L_X/8+\gamma(L_Y-L_X)/16\), and the explicit shoulders cost
at most \(0.88(L_Y-L_X)\). The denominator contributes
\(-(5\gamma/24)L_Y+o(L)\). Their sum is
\[
 -\frac{16}{9}L_X+
 \left(0.88+\frac{\gamma}{16}-\frac{5\gamma}{24}\right)(L_Y-L_X)
 +o(L).
\]
Since \(\gamma=4/3\), \(\gamma L_X=\log s\), and
\(\gamma L_Y=\log b\), this is the claimed marked-pair exponent.
To remove the error coefficients, first make the reserve and smoothing
loss arbitrarily small, then make the buffer fraction sufficiently small,
and finally improve the core projection. Each choice remains fixed while
\(L\to\infty\). The regular and rare subsequence arguments cover every
sequence of site mixtures and wall ratios, proving uniformity.

\end{proof}

\begin{proof}[Proof of Theorem~\ref{thm:polygon-moment}]
The coefficient \(617/1200\) is smaller than \(2/3\). For the configurations selected in
Lemma~\ref{lem:07_magnetic_analysis-marks}, $b\asymp D$ and $s\le b$.
The sharp theorem therefore bounds the retained port-pair weights by
\(D^{o(1)}s^{-4/3}(D/s)^{2/3}\). At dyadic diameter \(D\), a dyadic
displacement scale \(s\) has only \(O(s^2)\) possible placements after
fixing the first port modulo translation. Its contribution is at most
\[
 D^{o(1)}s^2s^{-4/3}(D/s)^{2/3}=D^{2/3+o(1)}.
\]
There are \(O(\log D)\) displacement scales. The short-polygon contribution
is already bounded in Lemma~\ref{lem:07_magnetic_analysis-marks} by the
planar diameter tail. Summing the dyadic diameter bounds up to \(H\)
proves Theorem~\ref{thm:polygon-moment}.
\end{proof}
